% The master paper. Every venue cut (paper/CUTS.md) is a variant of THIS file that
% \input s a subset -- never a second draft. Overleaf is the compiler; this repo is the
% source of truth and is never edited from Overleaf.
% TMLR's official style (tmlr.sty/tmlr.bst/fancyhdr.sty, Apache-2.0, TMLR_STYLE_LICENSE,
% from github.com/JmlrOrg/tmlr-style-file). With no option it is the ANONYMOUS submission:
% the author block below is hidden. [preprint] de-anonymises for arXiv; [accepted] is the
% camera-ready. tmlr.sty sets the page (6.5in text width) and loads natbib itself.
\documentclass[10pt]{article}
\usepackage{tmlr}
% \usepackage[preprint]{tmlr}

\usepackage{amsmath,amssymb,amsthm}
\usepackage{booktabs}
\usepackage{array}
\newcolumntype{L}[1]{>{\raggedright\arraybackslash}p{#1}}
\usepackage{graphicx}
\usepackage{url}
\usepackage[colorlinks=true,allcolors=blue]{hyperref}

\newtheorem{theorem}{Theorem}
\newtheorem{proposition}[theorem]{Proposition}
\newtheorem{lemma}[theorem]{Lemma}
\newtheorem{corollary}[theorem]{Corollary}
\theoremstyle{definition}
\newtheorem{definition}{Definition}
\theoremstyle{plain}

% Table glyphs, so the modulus table's cells stay bare (register sec. 6).
\newcommand{\yes}{\checkmark}
\newcommand{\no}{--}
\DeclareMathOperator{\Jclass}{J}

% "Causal" was dropped from this title 2026-09-21. The ablation FFTs the SAVED LOGIT
% TENSOR -- it re-runs no forward pass, ablates no weight, patches no activation -- so the
% word was doing work the protocol does not do. The evidence class is stated precisely in
% the abstract and in Section 5 instead of asserted in three words here.
\title{A Clock on Every Stratum:\\
Character Circuits for Modular Multiplication\\
at Non-Square-Free Moduli}
% Hidden in the submission build (tmlr with no option). Single author: the researcher
% decided 2026-09-26 that there is no supervisor co-author for now.
% The author block is completed at submission; this repository
% carries no authorship metadata by design. See README.
\author{}

\begin{document}
\maketitle

\begin{abstract}
% STE descriptive
Mechanistic accounts of modular arithmetic in transformers largely stop where global
invertibility stops. \citet{chen2026multiplication} extend the character account to
$a \cdot b \bmod n$. They treat $\mathbb{Z}/n\mathbb{Z}$ as a monoid, and its Green's
$\mathcal{J}$-classes stratify the input space. Their evidence covers four square-free moduli,
and they describe it as only correlational. Non-square-free moduli have non-regular classes, and
these classes contain nilpotents. \citet{chen2026multiplication} name these moduli as the case
that breaks the reduction to local group characters.

We show that the reduction does not break. With $x = d \cdot u$, a $\mathcal{J}$-class is a
torsor under the \emph{local} units $(\mathbb{Z}/(n/d)\mathbb{Z})^{\times}$. Under the global
units the action is transitive but not free. So every stratum is multiplication in a smaller
ring, carried on that stratum's own local group $G_{m} = (\mathbb{Z}/(n/m)\mathbb{Z})^{*}$.
Nilpotents collapse a stratum onto a smaller local group. They do not destroy its character structure.

We measure this across $23$ moduli, $14$ of them non-square-free and $6$ of them prime powers.
A transformer runs the same character-based clock on every stratum. The measurement covers $32$
strata, $157$ stratum-run measurements and $12$ distinct local groups. We committed its
criteria before any per-stratum number existed. At $n = 165$ the strata raise the share of the
multiplication table that the mechanism accounts for from $23.5\%$ to $82.3\%$.

It is also the \emph{same} clock, and not one clock per stratum. Strata that share a local group
select exactly the same characters of it in $158$ of $160$ pre-registered pairs. Against this, a
$10{,}000$-draw null is at its floor, and a control that destroys the structure reads $0$ of
$160$.

The output depends on those frequencies, and they are not incidental. A $10{,}000$-draw
permutation ablation of the logit spectrum puts the key set beyond \emph{every} draw at $133$ of
$157$ stratum-measurements. It gives $\hat{p} \le 0.0023$ at all $157$, with excluded/baseline
from $3.98 \times 10^{2}$ to $2.71 \times 10^{8}$. It also holds in $28$ of $29$ further
single-modulus runs, and we name the one exception. It reaches $2^{7}$, the deepest nilpotency
tested, in product-character coordinates where no discrete log exists.

This is an intervention on the saved logit spectrum, the restricted/excluded protocol of
\citet{nanda2023progress}. It re-runs no forward pass, ablates no weight and patches no
activation. So it establishes what the output depends on, and not which internal component
computes it. We therefore claim no more than that protocol earns, and we name it wherever we
cite the evidence.

At two moduli we do claim more: we delete those characters from the \emph{embedding} and run
the network forward unmodified. This takes unit-grid accuracy to $0.0089$--$0.0091$ at the prime
$n = 113$ (chance $1/112$). It takes it to $0.0000$--$0.0120$ at $n = 121 = 11^{2}$ (chance
$1/110$), and the edit leaves the non-regular stratum of that modulus intact. Both results hold
on three of three seeds at each modulus. The complementary edit keeps only those characters and
leaves accuracy at $0.9596$ or above. This is a weight-space intervention with both
nonlinearities downstream of the edit, on checkpoints that reproduce the archived runs
bit-identically.

The algebra of $n$ changes how many strata there are and how they nest, and not the mechanism
that runs on them. Two negatives bound the claim. The criterion that our own pre-registration
named primary does not discriminate a grokked model from a failed one. So the verdict rests on
the three criteria that do. The clearest early signal that we found also rises identically in
runs that never generalise. It is therefore not a progress measure.
% STE end
\end{abstract}

\section{Introduction}
\label{sec:intro}

Grokking, where generalisation arrives long after a training set is fit, is useful here as an
instrument rather than as a subject. A small transformer trained to convergence on a closed
algebraic task can be reverse-engineered completely, and the resulting circuit is a statement
about what the architecture does with the structure it is given
\citep{power2022grokking,nanda2023progress}.

Modular multiplication is the harder case, and a change of basis is the source. For
$a + b \bmod n$ the network represents inputs by characters of $\mathbb{Z}/n\mathbb{Z}$ and
adds phases. For $a \cdot b \bmod n$ the same circuit exists but lives in the multiplicative
character basis, reached by discrete log, and a model that is sharply sparse there looks dense
in raw residue coordinates \citep{nguyen2026dlogclock}. The basis is not a presentational
choice: a statistic read in the wrong one inverts the conclusion.

That account needs the units to form a group, and for a general modulus most of the
multiplication table is not units. \citet{chen2026multiplication} took the step past it: they
model $(\mathbb{Z}/n\mathbb{Z}, \cdot)$ as a monoid, show that character-style computation
localises to disjoint algebraic strata identified with Green's $\mathcal{J}$-classes, and
demonstrate class-sensitive routing. That framing is theirs. Independent derivation is
not claimed and would not be a defence: \citet{chen2026multiplication} published the stratified
framing before this work reached it, and they are cited as its origin. Their result covers four
moduli, all square-free, and they state its two limitations themselves: the evidence is
\emph{only correlational}, and non-square-free moduli introduce non-regular
$\mathcal{J}$-classes containing nilpotent elements, which they describe as breaking the
reduction to local group characters. This paper answers both. It is the non-square-free
completion of their programme, and it replaces correlation with an ablation: the first
of the two instruments they ask for, not the second (Sections~\ref{sec:causal} and
\ref{sec:strata}).

The reduction does not break. Writing $x = d \cdot u$ makes a
$\mathcal{J}$-class a torsor under the \emph{local} units, so a stratum is multiplication
in a smaller ring on that stratum's own local group $G_{m}$. Nilpotents \emph{collapse} a stratum onto a
smaller local group; they do not destroy the character structure
(Theorem~\ref{thm:stratum}, Section~\ref{sec:strata}). The algebra of $n$
does not change the mechanism. It changes how many strata there are and how they nest: one
for a field, a nested $p$-adic chain for a prime power, a divisor lattice for a square-free
composite.

% STE procedural
\begin{figure}[t]
  \centering
  \includegraphics[width=\textwidth]{../figures/F0_pipeline.pdf}
  \caption[What this paper does.]{\textbf{What this paper does.} Train a one-layer transformer to grok
    $a \cdot b \bmod n$ at each of $23$ moduli. Then read the trained model in the coordinate
    the algebra supplies. The first three steps are standard. The fourth is not. Relabel each
    unit by its exponent tuple. Only then does multiplication of units become addition.
    Theorem~\ref{thm:stratum} gives every stratum $J_{d} \times J_{e}$ its own such
    coordinate, on its own local group $G_{m}$.\par
    Measure and ablate per stratum, in that stratum's coordinate and not in the raw residues
    of $\mathbb{Z}/n\mathbb{Z}$. The figure regenerates from
    \texttt{src/viz/plots.py::pipeline}.}
  \label{fig:pipeline}
\end{figure}
% STE end

\paragraph{Contributions.}
\begin{enumerate}
  \item \emph{The stratum identity} (Theorem~\ref{thm:stratum}), proved and also
    checked by exhaustion, which
    turns ``non-regular classes break the reduction'' into ``non-regular classes carry the
    reduction on a smaller group'', and fixes what to measure without any freedom of choice.
  \item \emph{The mechanism, measured on every stratum}: $32$ strata across six moduli, $157$
    stratum-run measurements and $12$ distinct local groups, with criteria committed before any
    per-stratum number existed. It raises the fraction of the multiplication table accounted
    for from $23.5\%$ to $82.3\%$ at $n = 165$ and from $64.0\%$ to $89.6\%$ at $n = 125$
    (Section~\ref{sec:strata}, five seeds).
  \item \emph{Ablation evidence across $23$ moduli, $14$ of them non-square-free}, including
    all six prime powers (Section~\ref{sec:causal}, five seeds at six moduli, three at the
    rest). It reaches $2^{7}$, nilpotency depth $7$, where the key set is beyond all
    $10{,}000$ draws in $2/2$ runs (improvement $2.77$--$476\times$) in product-character
    coordinates where no discrete log exists.
  \item \emph{One mechanism, not one per stratum}: strata that share a local group select
    exactly the same characters of it, $158$ of $160$ pairs across five seeds, with a
    cell-shuffle control at $0$ of $160$. This is the test that separates a single indexed
    mechanism from a private sub-circuit per block (Section~\ref{sec:strata-sameness}). It
    ships without a failed-run control, and Section~\ref{sec:strata-sameness} says why.
  \item \emph{The negative results, reported as results}. The pre-registered primary
    criterion of our own stratum test does not discriminate a grokked model from a failed one,
    and the verdict rests on the three criteria that do. The early CRT enrichment ramp is not a
    progress measure. Five earlier claims of ours are retracted in
    Appendix~\ref{app:retractions}, with what killed each and the control that would have
    caught it.
\end{enumerate}

Section~\ref{sec:related} relates this work to prior mechanistic accounts.
Section~\ref{sec:setup} defines the task and its strata, and Section~\ref{sec:methods} the
model, the coordinates, the statistics and the ablation protocol. Section~\ref{sec:calib}
validates the pipeline against published results and compares our two implementations.
Section~\ref{sec:clock} reports where models grok and the clock at prime powers,
Section~\ref{sec:causal} the ablation evidence, Section~\ref{sec:strata} the measurement on
every stratum and Section~\ref{sec:crt} the additive CRT structure.
Section~\ref{sec:discussion} discusses the results, Section~\ref{sec:limits} states the
limitations and Section~\ref{sec:conclusion} concludes.

\section{Related work}
\label{sec:related}

% STE descriptive
Table~\ref{tab:methods} compares the seven studies closest to this one. It gives the task, the
algebra, the moduli, the coordinates of the analysis and the kind of evidence that each study
reports.

\begin{table}[t]
  \centering
  \footnotesize
  \caption{\textbf{What each method can see.} Each row is one study, and every cell comes
    from that study's PDF. The evidence column gives the tests that the study reports. A null
    arm is a control condition for which the account predicts no effect. ``Not stated'' means
    that the PDF does not contain the word pre-registration.}
  \label{tab:methods}
  \setlength{\tabcolsep}{3pt}
  \renewcommand{\arraystretch}{1.15}
  \begin{tabular}{L{0.10\textwidth} L{0.055\textwidth} L{0.12\textwidth} L{0.14\textwidth}
      L{0.13\textwidth} L{0.15\textwidth} L{0.12\textwidth} L{0.075\textwidth}}
    \toprule
    study & task & algebra & moduli (square-free?) & coordinates & evidence & null arm & pre-reg. \\
    \midrule
    \citet{nanda2023progress} & $a + b$ & $\mathbb{Z}/p\mathbb{Z}$ under $+$ & $113$, $53$, $109$ (primes) & Fourier basis & logit ablation; activation replacement & non-key frequencies & not stated \\
    \citet{zhong2023clock} & $a + b$ & $\mathbb{Z}/p\mathbb{Z}$ under $+$ & $59$ (prime) & Fourier basis; principal components & correlational; rank-2 embedding replacement & none reported & not stated \\
    \citet{chughtai2023toy} & $a b$ in a group & finite groups & $C_{113}$, $C_{118}$ (square-free), and five non-cyclic groups & irreducible representations & logit ablation; activation replacement & directions predicted to be noise & not stated \\
    \citet{nguyen2026dlogclock} & $a \cdot b$ & $(\mathbb{Z}/p\mathbb{Z})^{*}$ & $113$; ten primes from $59$ to $113$ & discrete-log characters & correlational & none reported & not stated \\
    \citet{chen2026multiplication} & $a \cdot b$ & the monoid $(\mathbb{Z}/n\mathbb{Z}, \cdot)$ & $113$, $143$, $154$, $165$ (all square-free) & characters of each $\mathcal{J}$-class group & correlational, in their words & random vocabulary subsets & not stated \\
    \citet{hwang2026pfe} & $a + b$ & $\mathbb{Z}/N\mathbb{Z}$ under $+$, by CRT & primes; ten composites from $15$ to $385$ (all square-free) & prime-indexed Fourier features, fixed & input-channel ablation, MLP classifier & irrelevant channels; shuffled features & not stated \\
    this work & $a \cdot b$ & the monoid $(\mathbb{Z}/n\mathbb{Z}, \cdot)$ & $23$ moduli, $14$ not square-free & characters of each stratum's local group $G_{m}$ & logit ablation; weight intervention at $n = 113$, $121$ & random character sets & $9$ of the $14$ rows of Table~\ref{tab:glance} \\
    \bottomrule
  \end{tabular}
\end{table}
% STE end

\paragraph{Grokking as an instrument.} \citet{power2022grokking} established the phenomenon:
on small algorithmic datasets, generalisation can appear long after the training set is fit,
and the smaller the training fraction the longer the delay. That paper claims a testbed rather
than a mechanism, and the testbed is what the literature since has used.
\citet{nanda2023progress} supplied the mechanism for $a + b \bmod n$ by reverse-engineering a
one-layer transformer completely: the network embeds inputs in Fourier features, adds phases
through the attention and MLP, and reads out with a trigonometric identity, confirmed by
ablations in Fourier space. Their restricted and excluded losses are the ablation protocol we
use, in a coordinate they did not need.

\paragraph{The algorithm is not inevitable.} \citet{zhong2023clock} showed that the Clock
algorithm is not the only one a network may find on the same task: a Pizza algorithm is
prevalent under other hyperparameters, algorithmic phase transitions exist, and networks
sometimes ensemble. \citet{chughtai2023toy} generalised the setting to arbitrary finite groups
and found universality to be two-level: the circuit \emph{family} is characterisable by
representation theory, while the particular circuit and the order in which it develops are
arbitrary. Their paired-ablation design, testing both the signal the theory predicts and the
noise it predicts, is the standard we hold our own ablations to.
\citet{stander2024cosets} analyse the same territory through cosets.

\paragraph{Multiplication and the basis.} Multiplication modulo $n$ resisted this treatment
for a simple reason: in raw residue coordinates a grokked multiplicative model has a dense
spectrum. That observation is \citet{furuta2024towards}'s, whose Section~5.2 reports that
multiplication ``leverages all frequencies'' with cosine-biased components, against the sparse
key-frequency picture addition gives.
\citet{nguyen2026dlogclock} showed the density is a basis artefact. Re-analysed
in the multiplicative character basis, where discrete log turns multiplication into addition,
the same model is sparse, and they report $\mathrm{Gini}\ 0.58$ against $0.07$, four key
frequencies, and $96.9\%$ of neurons single-frequency. Section~\ref{sec:calib} replicates
their spectral statistics and withdraws one clause of our own earlier replication.
\citet{doshi2024grokking} classify which modular polynomials are learnable and offer it as a
hypothesis with supporting evidence; their inverse participation ratio is a third convention,
reciprocal and bounded in $[0,1]$, which is why Section~\ref{sec:methods-sparsity} states ours
explicitly. \citet{furuta2024towards} and \citet{kunin2024getrich} supply the initialisation
and lazy-to-rich picture our initialisation control follows.
The lesson generalises past spectra, and we state it because the cost is easy to underestimate:
\emph{any} diagnostic calibrated on addition can read a multiplicative model as structureless
while the structure is present in exponent coordinates. We hit this ourselves. The same
grokked model's top eight logit components read $17.0\%$ ``$a+b$'' raw and $98.4\%$ in discrete-log coordinates
(\S\ref{sec:methods-sparsity}), and we record it as a methodological hazard rather than a
finding about multiplication.

\paragraph{Divisor structure is something these models already organise around.}
Before any character account, \citet{charton2024learning} trained transformers to compute
$\gcd(a, b)$ and characterised their predictions completely. The model learns a set $D$ of
integers and returns the largest element of $D$ dividing both inputs, so that every pair
sharing a $\gcd$ receives the same output, \emph{including where that output is wrong}.
The equivalence structure therefore sits in the computation rather than only in the target.
That is the same shape as Consequence~2 of Theorem~\ref{thm:stratum}, reached on a different
task, in base-$10$ digit coordinates, and with no representation-theoretic model; our fibre
residual \eqref{eq:resid} is the quantitative version of the constancy he observes
qualitatively. We read it as independent evidence that divisor classes are a structure these
models find, and note that he also reports small primes being grokked only after extended
training.

\paragraph{The stratified framing, which is not ours.}
\citet{chen2026multiplication} extend the character account past the group case. Where global
invertibility fails they model $(\mathbb{Z}/n\mathbb{Z}, \cdot)$ as a monoid, show that
group-character-style computation \emph{localises to disjoint algebraic strata that partition
the input space}, identify those strata with Green's $\mathcal{J}$-classes, and demonstrate
class-sensitive attention routing. They do so on four moduli, $113$, $143$, $154$ and $165$,
every one of which their Table~1 marks square-free. Their Theorem~3.4 requires each class to \emph{be} a group,
\begin{equation}
  \label{eq:chen}
  J_{d} \;\cong\; (\mathbb{Z}/(n/d)\mathbb{Z})^{\times}
  \quad\text{with a local identity } e_{J_{d}} \text{ and a local inverse } c^{\sharp},
\end{equation}
and their account of the read-out, developed in their Section~4 rather than in
Theorem~3.4, evaluates $\mathrm{Logit}(c) \propto \chi_{\rho_{d}}(a\,b\,c^{\sharp})$ on that
group. A non-regular class supplies neither: with no idempotent there is no $c^{\sharp}$.
Proposition~\ref{prop:torsor} replaces the isomorphism in \eqref{eq:chen} by a
\emph{bijection} that needs no group structure, and Theorem~\ref{thm:stratum} recovers the
character account on $G_{m}$ instead of on $J_{d}$. This is not a counterexample to Clifford--Munn--Ponizovski\u{\i}. CMP
\citep{steinberg2016monoids}, which \citet{chen2026multiplication} state as their
Theorem~D.16, says that the irreducible representations of a finite monoid are indexed by the
\emph{regular} $\mathcal{J}$-classes: a non-regular class contributes none, and no argument
here disputes that. What Theorem~\ref{thm:stratum} factors through $G_{m}$ is the
\emph{function the network has to compute} on a stratum, not the monoid's representation
theory. The characters we measure are characters of the local group $G_{m}$, which is a
genuine group; they are not irreps of $(\mathbb{Z}/n\mathbb{Z}, \cdot)$ attached to a
non-regular class, and the reduction does not require any. What is new here is not a stronger theorem but a wider one, and the width is the
contribution. Proposition~\ref{prop:torsor} trades a group isomorphism for a bijection of
sets: weaker hypotheses and a weaker conclusion, which is a different statement and not a
strengthening of their Theorem~3.4, and we do not claim it as one. Two things follow that
their theorem does not give. \emph{First, non-vacuity exactly where theirs is vacuous.}
\eqref{eq:chen} asks $J_{d}$ to be a group, so at a non-square-free modulus it is silent on
precisely the classes that make such a modulus interesting: $8$ of the $16$ at $n = 120$,
and the ones containing the nilpotents. \eqref{eq:torsor} and Theorem~\ref{thm:stratum}
are stated for every $d \mid n$ and ask nothing beyond $\mathcal{J}$-classhood. \emph{Second, cross-stratum composition.}
Theorem~\ref{thm:stratum} is a statement about the pair $J_{d} \times J_{e}$ for all
$d, e$, which their Theorem~3.4 does not address even in the square-free case. It says which
class a block lands in, \eqref{eq:jcompose}, and which local group carries the
multiplication once it is there, \eqref{eq:stratum}. That is the statement our
measurements use, because the strata of the input grid are \emph{blocks} and not
classes, and every per-stratum number in Section~\ref{sec:strata} is measured on a block.
The framing is theirs, and this
paper is its completion, not its rediscovery. They state two limitations of their own result in their own
words: that their evidence is only correlational, and that non-square-free moduli introduce
non-regular $\mathcal{J}$-classes containing nilpotent elements, which they describe as
breaking the reduction to local group characters. Sections~\ref{sec:causal} and
\ref{sec:strata} supply the two.

The relation to \citet{chughtai2023toy} is worth naming precisely, because it is the shape of
our result. They argue universality is two-level: the circuit family invariant, the specific
circuit arbitrary. What we add is that for $a \cdot b \bmod n$ the two levels are
\emph{algebraic}: the family is the character clock, and the index set is fixed by the
stratum's own local group.

\paragraph{The CRT connection is not ours either.} \citet{hwang2026pfe} supply a
prime-indexed embedding and prove block-diagonality by Schur's lemma, with an ablation carrying
a null arm: task-relevant channels drop $0.60$--$0.92$ while irrelevant channels sit at or
below the noise floor. Theirs is the basis \emph{supplied}. What Section~\ref{sec:crt-law}
adds is that the equivalent structure is \emph{learned} rather than imposed, on
\emph{multiplication} rather than addition, and at \emph{non-square-free} moduli, where the
predicted set is built from maximal prime powers rather than primes and the distinction is
testable.

\section{Setup}
\label{sec:setup}

\subsection{Task}
\label{sec:setup-task}

We train a transformer on the multiplication table of $\mathbb{Z}/n\mathbb{Z}$. Each example
is the token sequence $(a, b, {=})$ with $a, b \in \mathbb{Z}/n\mathbb{Z}$, and the model
\begin{equation}
  \label{eq:task}
  f_{\theta} : (\mathbb{Z}/n\mathbb{Z})^{2} \longrightarrow \mathbb{R}^{n},
  \qquad
  \hat{y}(a,b) \;=\; \arg\max_{c} \; f_{\theta}(a,b)_{c},
  \qquad
  y(a,b) \;=\; a \cdot b \bmod n ,
\end{equation}
is trained by cross-entropy $\mathcal{L}(\theta) = -\frac{1}{|T|}\sum_{(a,b) \in T}
\log \mathrm{softmax}\bigl(f_{\theta}(a,b)\bigr)_{y(a,b)}$ on a training set $T$, with the
logits read off the final position. The vocabulary is $n+1$ symbols and the
output layer has $n$ classes. All $n^{2}$ pairs are enumerated, a fixed random $30\%$ is held
out for training and the remaining $70\%$ is the test set, so generalisation is measured on
the majority of the table.

The operation is the whole reason for the setup. Addition modulo $n$ is a group operation:
every element is invertible, and a network that represents $a$ and $b$ by characters of
$\mathbb{Z}/n\mathbb{Z}$ can compute $a+b$ by adding phases
\citep{nanda2023progress}. Multiplication is not a group operation unless $n$ is prime.
$(\mathbb{Z}/n\mathbb{Z}, \cdot)$ is a commutative monoid; its invertible elements are the
units
\begin{equation}
  \label{eq:units}
  (\mathbb{Z}/n\mathbb{Z})^{\times} \;=\; \{x : \gcd(x, n) = 1\},
  \qquad
  \bigl| (\mathbb{Z}/n\mathbb{Z})^{\times} \bigr| \;=\; \varphi(n),
\end{equation}
and every other element is a zero divisor. On the units
the same phase-addition argument applies in exponent coordinates, and the resulting circuit is
the discrete-log clock \citep{nguyen2026dlogclock}. Off the units it does not apply at all,
because there are no local inverses to build characters from. What the network does on the
$n^{2} - \varphi(n)^{2}$ non-unit pairs is the question this paper is about.

\subsection{Strata}
\label{sec:setup-strata}

\begin{definition}[Green's $\mathcal{J}$-classes and strata]
\label{def:jclass}
Green's relations are the standard structural decomposition of a semigroup
\citep{howie1995semigroups}. For $d \mid n$ put
\begin{equation}
  \label{eq:jclass}
  J_{d} \;=\; \{x \in \mathbb{Z}/n\mathbb{Z} : \gcd(x, n) = d\},
  \qquad
  |J_{d}| \;=\; \varphi(n/d),
  \qquad
  \mathbb{Z}/n\mathbb{Z} \;=\; \bigsqcup_{d \mid n} J_{d} .
\end{equation}
These are the $\mathcal{J}$-classes of the monoid $(\mathbb{Z}/n\mathbb{Z}, \cdot)$, and they
compose \emph{exactly}:
\begin{equation}
  \label{eq:jcompose}
  J_{d} \cdot J_{e} \;=\; J_{\gcd(de,\, n)} .
\end{equation}
Containment is Theorem~\ref{thm:stratum}: it writes $xy = m\,(w\,uv \bmod q)$ with
$m = \gcd(de,n)$ and $w\,uv$ a unit modulo $q = n/m$, so $\gcd(xy, n) = m$ for every
$x \in J_{d}$ and $y \in J_{e}$. Equality then needs only surjectivity: $q$ divides $n/d$,
so reduction carries $(\mathbb{Z}/(n/d)\mathbb{Z})^{\times}$ onto
$(\mathbb{Z}/q\mathbb{Z})^{\times}$, and taking $v = 1$ lets $w\,uv$ range over every unit
modulo $q$. Verified exhaustively at every divisor pair of every $n \le 200$
($8{,}253$ pairs, $0$ mismatches).
The induced partition of the $n \times n$ input grid into the blocks
$J_{d} \times J_{e}$ is what we call its \emph{strata}. A class is \emph{regular} when it
contains an idempotent, $\exists\, \varepsilon \in J_{d}$ with $\varepsilon^{2} = \varepsilon$,
the condition under which local inverses, and hence local group characters, exist.
\end{definition}

\citet{chen2026multiplication}
prove (their Proposition D.13) that every $\mathcal{J}$-class is regular when $n$ is square-free,
and state the converse case as the limitation of their work. Non-square-free moduli admit
non-regular classes containing \emph{nilpotent} elements ($x$ with $x^{k} = 0$ for some
$k \ge 1$, which exist in $\mathbb{Z}/n\mathbb{Z}$ exactly when $n$ is not square-free), and
they describe these as breaking the reduction to local group characters. Our implementation reproduces their theorem on their own
moduli $143$, $154$ and $165$, and finds non-regular classes at exactly the moduli where the
theorem permits them: the counts for $113, 119, 121, 125, 120$ are $0, 0, 1, 2, 8$
(\texttt{src/tasks/algebra.py::j\_structure}). Across all $23$ moduli we train, every
square-free modulus has zero non-regular classes and every non-square-free modulus has at
least one.

The difference the composition rule makes is visible in one line. For square-free $n$ the
diagonal is closed and each stratum is a group; for a prime power it descends toward zero:
\begin{equation}
  \label{eq:descent}
  \underbrace{J_{d} \cdot J_{d} = J_{d}}_{n \text{ square-free}}
  \qquad\text{versus}\qquad
  \underbrace{J_{11} \cdot J_{11} = \{0\}}_{n = 121}
  \qquad
  \underbrace{J_{5} \cdot J_{5} \to J_{25} \to \{0\}}_{n = 125} .
\end{equation}
Writing $\nu(n) = \min\{k \ge 1 : x^{k} = 0 \text{ for every nilpotent } x \in
\mathbb{Z}/n\mathbb{Z}\}$ for the nilpotency depth, $\nu(121) = 2$ and $\nu(125) = 3$. The two
moduli differ in $\nu$ at nearly identical $\varphi(n)$ and identical model and dataset size,
which makes them a graded pair rather than two arbitrary composites.

% STE descriptive
\begin{figure}[t]
  \centering
  \includegraphics[width=\textwidth]{../figures/F5_descent.pdf}
  \caption[The square map walks a class down the divisor lattice.]{\textbf{The square map walks a class down the divisor lattice.} Each node is a
    $\mathcal{J}$-class, placed by the number of prime factors of $d$ counted with
    multiplicity. Each edge is the map $J_{d} \mapsto J_{\gcd(d^{2},\, n)}$ of
    \eqref{eq:jcompose}. A teal loop marks a class that squares to itself, which is
    \eqref{eq:descent}'s square-free case, and a grey edge marks one that does not. Node fill
    and outline give regularity.\par
    At $n = 121$ the single non-regular class reaches $\{0\}$ in one step, and at $n = 125$
    it takes two. That is the difference between $\nu(121) = 2$ and $\nu(125) = 3$ at
    $\varphi(121) = 110$ against $\varphi(125) = 100$. At $n = 120$ the lattice is wide
    rather than deep: $16$ classes, $8$ of them non-regular. The figure regenerates from
    \texttt{src/viz/plots.py::descent\_lattice}, computed from
    \texttt{src/tasks/algebra.py}.}
  \label{fig:descent}
\end{figure}
% STE end

% STE descriptive
\begin{figure}[t]
  \centering
  \includegraphics[width=\textwidth]{../figures/F2_stratification.pdf}
  \caption[The same operation over five algebras.]{\textbf{The same operation over five algebras.} Each panel is the
    multiplication table of $\mathbb{Z}/n\mathbb{Z}$ for one of our moduli. Rows and columns
    are ordered by $\gcd(x, n)$, so each stratum $J_{d} \times J_{e}$ is a rectangle, and
    cell shade is the $\mathcal{J}$-class of the product. The operation, the architecture and
    the training budget are identical across the five. So only the algebra of $n$ differs.
    There is one stratum at a prime, $9$ at $11^{2}$, $16$ at $5^{3}$ and at $7 \cdot 17$,
    and $256$ at $2^{3} \cdot 3 \cdot 5$.\par
    Red outlines mark non-regular classes, which occur at exactly the three non-square-free
    moduli and at neither square-free one. That is Proposition D.13 of
    \citet{chen2026multiplication}, drawn. The figure regenerates from
    \texttt{src/viz/plots.py::stratification}, computed from \texttt{src/tasks/algebra.py}.}
  \label{fig:stratification}
\end{figure}
% STE end

\subsection{A stratum is multiplication in a smaller ring}
\label{sec:setup-theorem}

The reduction to local characters does not in fact break. It relocates, and two statements
say where to: first the indexing that makes a non-regular class usable, then the identity that
says what multiplication does in it.

\begin{proposition}[$\mathcal{J}$-classes are torsors under the \emph{local} units]
\label{prop:torsor}
For every $d \mid n$ the map
\begin{equation}
  \label{eq:torsor}
  \iota_{d} : (\mathbb{Z}/(n/d)\mathbb{Z})^{\times} \longrightarrow J_{d},
  \qquad u \longmapsto d\,u \bmod n ,
\end{equation}
is a bijection, and it is defined whether or not $J_{d}$ is regular.
\end{proposition}

\begin{proof}
Since $d \mid n$ we have $\gcd(du, n) = d \cdot \gcd(u, n/d)$, which is $d$ exactly when
$\gcd(u, n/d) = 1$; so $\iota_{d}$ lands in $J_{d}$. Conversely any $x \in J_{d}$ is
$d\,(x/d)$ with $\gcd(x/d, n/d) = 1$. The two constructions are mutually inverse modulo
$n/d$, and counting gives $|J_{d}| = \varphi(n/d)$, which is \eqref{eq:jclass}.
\end{proof}

Proposition~\ref{prop:torsor} asks nothing of $J_{d}$ beyond being a $\mathcal{J}$-class. A
regular class is a group and carries its own characters. A non-regular class is not a group
and has none, but it is still a set on which the \emph{local} units
$(\mathbb{Z}/(n/d)\mathbb{Z})^{\times}$ act simply transitively, and
\eqref{eq:torsor} is that action written as a coordinate. The theorem says what
multiplication looks like in that coordinate.

\begin{theorem}[Stratum identity]
\label{thm:stratum}
Let $x \in J_{d}$ and $y \in J_{e}$, and write $x = d\,u$, $y = e\,v$ as in
\eqref{eq:torsor}. Put
\begin{equation}
  \label{eq:mw}
  m \;=\; \gcd(de,\, n), \qquad de \;=\; m\,w, \qquad q \;=\; n/m .
\end{equation}
Then $d \mid m$, $\;e \mid m$, $\;\gcd(w, q) = 1$, $\;uv$ is a unit modulo $q$, and
\begin{equation}
  \label{eq:stratum}
  x \cdot y \bmod n \;=\; m \cdot \bigl( w \cdot (u v) \bmod q \bigr).
\end{equation}
\end{theorem}

\begin{proof}
Fix a prime $p$ and write $\alpha = v_{p}(de)$, $\beta = v_{p}(n)$ for the $p$-adic
valuations. By \eqref{eq:mw}, $v_{p}(m) = \min(\alpha, \beta)$, hence
\begin{equation}
  \label{eq:valuations}
  v_{p}(w) \;=\; \max(\alpha - \beta,\, 0),
  \qquad
  v_{p}(q) \;=\; \max(\beta - \alpha,\, 0) .
\end{equation}
At most one of the two is positive, so no prime divides both $w$ and $q$, giving
$\gcd(w, q) = 1$.

Since $d \mid de$ and $d \mid n$ we have $d \mid \gcd(de, n) = m$, and likewise $e \mid m$;
therefore $q = n/m$ divides both $n/d$ and $n/e$. As $\gcd(u, n/d) = 1$ and
$\gcd(v, n/e) = 1$, reduction gives $\gcd(u, q) = \gcd(v, q) = 1$, so $uv$ is a unit
modulo $q$.

For the identity, $xy = (du)(ev) = m\,w\,uv$ as integers, and for any integer $t$ and any
$m \mid n$,
\begin{equation}
  \label{eq:mscale}
  (m t) \bmod n \;=\; (m t) \bmod (m q) \;=\; m \cdot (t \bmod q),
\end{equation}
since subtracting a multiple of $mq$ from $mt$ subtracts a multiple of $q$ from $t$. Applying
\eqref{eq:mscale} with $t = w\,uv$ gives \eqref{eq:stratum}.
\end{proof}

Three consequences follow, and every per-stratum measurement in this
paper is built on them.

\begin{enumerate}
  \item \emph{A stratum is multiplication in a smaller ring.} By \eqref{eq:stratum} the
    block $J_{d} \times J_{e}$ computes $(u,v) \mapsto w \cdot (uv)$ in
    $\mathbb{Z}/q\mathbb{Z}$, a fixed unit relabelling by $w$ of multiplication carried on the
    \emph{local group}
    \begin{equation}
      \label{eq:localgroup}
      G_{m} \;=\; (\mathbb{Z}/q\mathbb{Z})^{\times}, \qquad q = n/m, \qquad |G_{m}| = \varphi(q).
    \end{equation}
  \item \emph{The target is constant on fibres.} Reduction modulo $q$ gives a surjection
    \begin{equation}
      \label{eq:fibration}
      \pi_{d} : (\mathbb{Z}/(n/d)\mathbb{Z})^{\times} \twoheadrightarrow G_{m},
      \qquad
      y\bigl(\iota_{d}(u), \iota_{e}(v)\bigr) \;=\; m \cdot \bigl(w \, \pi_{d}(u)\pi_{e}(v)\bigr),
    \end{equation}
    so the target depends on $(u,v)$ only through $\bigl(\pi_{d}(u), \pi_{e}(v)\bigr)$. A
    stratum therefore has $|G_{m}|$ distinct targets and $\varphi(n/d)\,\varphi(n/e) \ge
    |G_{m}|$ cells, and the collapse is exactly this fibration.
  \item \emph{A clock must sit on the diagonal.} Writing $G_{m} \cong \mathbb{Z}_{o_{1}}
    \times \cdots \times \mathbb{Z}_{o_{r}}$ and indexing each unit by its exponent tuple,
    multiplication in $G_{m}$ is addition of tuples, so a character circuit on this stratum
    places its logit energy on the diagonal
    \begin{equation}
      \label{eq:diagonal}
      \Delta \;=\; \bigl\{ (\kappa, \kappa) : \kappa \in \widehat{G_{m}} \bigr\}
      \;\subset\; \widehat{G_{m}} \times \widehat{G_{m}},
      \qquad |\Delta| = |G_{m}| ,
    \end{equation}
    in that stratum's \emph{own} local-group coordinates, not in the raw residue
    coordinates of $\mathbb{Z}/n\mathbb{Z}$, where the components of $a+b$ and $a-b$ have no
    meaning on a multiplication task.
\end{enumerate}

$\Delta$ is fixed by the algebra before any data is seen, which is what lets
Section~\ref{sec:strata} ablate it without choosing it.

% STE descriptive
\begin{figure}[t]
  \centering
  \includegraphics[width=\textwidth]{../figures/F1_stratum.pdf}
  \caption[Theorem~\ref{thm:stratum}, and the limitation it answers.]{\textbf{Theorem~\ref{thm:stratum}, and the limitation it answers.}
    \emph{Top left:} at a square-free modulus every $\mathcal{J}$-class is regular, so
    $J_{3} \cdot J_{3} = J_{3}$ is closed and $J_{3}$ is a group, with identity $6$.
    \emph{Top right:} at $n = 20 = 2^{2} \cdot 5$ the class $J_{2}$ contains no idempotent.
    It is not a group and has no characters of its own, and $J_{2} \cdot J_{4} = J_{4}$
    leaves it. That is the case \citet{chen2026multiplication} exclude.\par \emph{Bottom:} the
    outlined stratum, unfolded. The substitution $x = 2u$ and $y = 4v$ of \eqref{eq:torsor}
    puts the same $16$ cells in torsor coordinates (panel 2). Each entry there is
    $w\,uv \bmod q$ with $m = 4$, $w = 2$ and $q = 5$. Multiplication by $m$ returns panel 1
    cell for cell, which is \eqref{eq:stratum}. The $16$ cells take $4$ distinct values,
    the order of the local group $G_{m} = (\mathbb{Z}/5\mathbb{Z})^{\times}$, and that
    collapse is the fibration \eqref{eq:fibration}. In $G_{m}$'s exponent coordinate
    multiplication is addition of $\kappa$ (panel 3), which is why a character circuit on
    this stratum sits on the diagonal \eqref{eq:diagonal}.\par
    Cell colour is the product and
    is identical in panels 1 and 2 because the identity holds cell by cell. The cell
    $x = 6$, $y = 12$ is ringed in all three. The figure regenerates from
    \texttt{src/viz/plots.py::stratum\_theorem}. Every value is computed from
    \texttt{src/tasks/algebra.py}, and the generator asserts \eqref{eq:stratum} and the
    non-regularity of $J_{2}$ on the block it draws.}
  \label{fig:stratum}
\end{figure}
% STE end

Nilpotents are what let $m = \gcd(de, n)$ rise to its ceiling $d\,e$ and drive
$n/m$ down; in the limit $n/m = 1$ and the stratum collapses onto the trivial group, which is
the ``one-way collapse'' of \citet{chen2026multiplication}. The collapse removes structure by
shrinking the local group. It does not replace character structure with something else.

Equation~\eqref{eq:stratum} is proved above. It is \emph{also} checked by exhaustion, because
a proof and an implementation can disagree. \texttt{algebra.py} verifies both clauses for
every one of the $n^{2}$ pairs at $n \in \{113, 121, 125, 120\}$ (the field, the two $p$-adic
chains and the widest divisor lattice in our set) as part of the repository's self-check
suite.

\subsection{Moduli}
\label{sec:setup-moduli}

Table~\ref{tab:moduli} is the set of moduli we train, and every later section refers back to
it. It extends Table~1 of \citet{chen2026multiplication} (modulus, prime decomposition,
$\omega(n)$, square-free) with the properties this paper separates results by: whether the
unit group is cyclic, $\varphi(n)$, zero-divisor density, the $\mathcal{J}$-class count and how
many of those classes are non-regular.

Of the $23$ moduli, $14$ are non-square-free and $6$ are prime powers, against the four
square-free moduli of \citet{chen2026multiplication}. The prime powers need two clauses rather
than one, because $n = 128 = 2^{7}$ is the only one whose unit group is non-cyclic:
$(\mathbb{Z}/128\mathbb{Z})^{*} \cong \mathbb{Z}_{2} \times \mathbb{Z}_{32}$, so it has no
discrete log, and a discrete-log statistic is not merely unmeasured there but undefined.
Results stated in discrete-log coordinates therefore cover the \emph{five cyclic} prime powers
$49, 81, 121, 125, 169$; the sixth is reached in exponent-tuple coordinates in
Section~\ref{sec:causal-128}.

% STE descriptive
\begin{table}[t]
  \centering
  \small
  \begin{tabular}{rlcccrrrrc}
    \toprule
    $n$ & factorisation & $\omega(n)$ & sq.-free & cyclic & $\varphi(n)$ & zdd &
      $|\mathcal{J}|$ & non-reg. & grokked \\
    \midrule
    49 & $7^{2}$ & 1 & \no & \yes & 42 & 0.143 & 3 & 1 & 0/3 \\
    54 & $2 \cdot 3^{3}$ & 2 & \no & \yes & 18 & 0.667 & 8 & 4 & 1/3 \\
    63 & $3^{2} \cdot 7$ & 2 & \no & \no & 36 & 0.429 & 6 & 2 & 1/3 \\
    75 & $3 \cdot 5^{2}$ & 2 & \no & \no & 40 & 0.467 & 6 & 2 & 2/3 \\
    81 & $3^{4}$ & 1 & \no & \yes & 54 & 0.333 & 5 & 3 & 3/3 \\
    91 & $7 \cdot 13$ & 2 & \yes & \no & 72 & 0.209 & 4 & 0 & 1/3 \\
    98 & $2 \cdot 7^{2}$ & 2 & \no & \yes & 42 & 0.571 & 6 & 2 & 3/3 \\
    99 & $3^{2} \cdot 11$ & 2 & \no & \no & 60 & 0.394 & 6 & 2 & 2/3 \\
    100 & $2^{2} \cdot 5^{2}$ & 2 & \no & \no & 40 & 0.600 & 9 & 5 & 2/3 \\
    105 & $3 \cdot 5 \cdot 7$ & 3 & \yes & \no & 48 & 0.543 & 8 & 0 & 3/3 \\
    113 & $113$ & 1 & \yes & \yes & 112 & 0.009 & 2 & 0 & 5/5 \\
    119 & $7 \cdot 17$ & 2 & \yes & \no & 96 & 0.193 & 4 & 0 & 5/5 \\
    120 & $2^{3} \cdot 3 \cdot 5$ & 3 & \no & \no & 32 & 0.733 & 16 & 8 & 5/5 \\
    121 & $11^{2}$ & 1 & \no & \yes & 110 & 0.091 & 3 & 1 & 5/5 \\
    123 & $3 \cdot 41$ & 2 & \yes & \no & 80 & 0.350 & 4 & 0 & 3/3 \\
    125 & $5^{3}$ & 1 & \no & \yes & 100 & 0.200 & 4 & 2 & 4/5 \\
    127 & $127$ & 1 & \yes & \yes & 126 & 0.008 & 2 & 0 & 3/3 \\
    128 & $2^{7}$ & 1 & \no & \no & 64 & 0.500 & 8 & 6 & 2/3 \\
    131 & $131$ & 1 & \yes & \yes & 130 & 0.008 & 2 & 0 & 3/3 \\
    143 & $11 \cdot 13$ & 2 & \yes & \no & 120 & 0.161 & 4 & 0 & 3/3 \\
    147 & $3 \cdot 7^{2}$ & 2 & \no & \no & 84 & 0.429 & 6 & 2 & 3/3 \\
    165 & $3 \cdot 5 \cdot 11$ & 3 & \yes & \no & 80 & 0.515 & 8 & 0 & 5/5 \\
    169 & $13^{2}$ & 1 & \no & \yes & 156 & 0.077 & 3 & 1 & 3/3 \\
    \bottomrule
  \end{tabular}
  \caption[The $23$ moduli.]{The $23$ moduli. \texttt{zdd} is zero-divisor density,
    $(n - \varphi(n))/n$, with $0$ counted. $|\mathcal{J}|$ is the number of Green's
    $\mathcal{J}$-classes, and \emph{non-reg.} is how many of those are non-regular. The
    latter is $0$ at every square-free modulus and positive at every non-square-free one, as
    \citet{chen2026multiplication} Proposition D.13 requires. \emph{grokked} counts seeds that
    reach test accuracy above $0.99$ under the single protocol of
    Section~\ref{sec:setup-training} ($30\%$ training fraction, $40{,}000$ steps).\par
    It excludes the arms that vary a hyperparameter. A run is classified by the median of its
    final ten logged samples, never by its last sample. Every column regenerates from
    \texttt{scripts/modulus\_table.py}.}
  \label{tab:moduli}
\end{table}
% STE end

The algebraic columns are computed, not tabulated: cyclicity is decided by an explicit
primitive-root search and $\mathcal{J}$-class regularity by searching each class for an
idempotent, and both are asserted against the design table in the repository's self-check
suite. Three claims in this project were lost to a confound between a grouping variable and a
size that travelled with it. The counts in Table~\ref{tab:moduli} ($\varphi(n)$, class size,
zero-divisor density) are exactly the sizes involved. A table that is regenerated rather than copied is the cheapest defence
available against the next one.

\section{Methods}
\label{sec:methods}

Three choices in this section are conventions rather than results, and each has cost this
project a wrong number at least once. They are the coordinate in which a spectrum is taken,
the power of the amplitude on which a sparsity statistic is computed, and what a random
control is allowed to randomise. We state each with its formula and its parameter, in the manner of
\citet{doshi2024grokking}'s footnote, because the literature does not agree on them.

\subsection{Model and training}
\label{sec:setup-training}

% STE descriptive
All runs use one transformer architecture from \citet{nanda2023progress}. The model has a
single layer, $d_{\text{model}} = 128$, four attention heads of width $32$, and a ReLU MLP of
width $512$. It has no layer normalisation or bias terms. The logits come from the final
position as $(z + \mathrm{ReLU}(z W_{\text{in}}) W_{\text{out}}) W_{U}$, where $z$ is the
post-attention residual stream. A learned positional embedding $W_{\text{pos}}$ adds to the
token embedding $W_{E}$, and attention uses a causal mask. Outside the sweep of
Section~\ref{sec:causal-init}, each weight matrix starts from a zero-mean Gaussian with
standard deviation $1/\sqrt{512}$ for $W_{\text{out}}$ and $1/\sqrt{128}$ for all other
matrices.

Training is full-batch AdamW with learning rate $10^{-3}$, weight decay $1.0$,
$\beta = (0.9, 0.98)$ and $\epsilon = 10^{-8}$. The loss is cross-entropy on the $30\%$
training split, for $40{,}000$ steps. The run seed sets both the initial weights and the
random split. PyTorch logs test accuracy and both losses every $200$ steps, and the engine
logs them every $100$ steps. Test accuracy is measured on the pairs held out from
training. The grokking step of a run is the first logged step at which test accuracy is above
$0.99$.

The weight decay is the setting under which grokking is reliably observed on this task
family \citep{power2022grokking}. It is held fixed across every modulus so that the modulus
is the only variable.

Two independent implementations produce the results in this paper. The breadth, $23$ moduli
with up to five seeds, comes from a PyTorch implementation. The same architecture is also
implemented on an autograd engine written from scratch for this project. The engine
reproduces the PyTorch forward pass and gradients to machine precision and is used as an
independent check at four moduli. Section~\ref{sec:impl} treats the two arms and the one
systematic difference between them.
% STE end

\subsection{Coordinates}
\label{sec:methods-coords}
\label{sec:setup-coords}

Two coordinate systems appear throughout, and every character index in this paper names which
one it is in. Where $(\mathbb{Z}/q\mathbb{Z})^{\times}$ is cyclic we fix a primitive root $g$
and use \emph{discrete-log coordinates}, $x = g^{\,\ell}$, in which multiplication is addition
of $\ell$. Where it is not cyclic, as at $n = 119$, $n = 120$, $n = 128$ and in every stratum
whose local group is a product, there is no discrete log. There we use \emph{exponent-tuple
coordinates}, given by a fixed decomposition into cyclic factors. The choice of $g$ is a choice
of labelling and not of structure; Proposition~\ref{prop:generator} records that
the readout is exactly invariant to it.

All spectral statistics on a multiplication run are taken in a coordinate in which
multiplication is addition. Where $(\mathbb{Z}/q\mathbb{Z})^{\times}$ is cyclic of order
$\varphi(q)$, a primitive root $g$ gives the discrete-log isomorphism
\begin{equation}
  \label{eq:dlog}
  \mathrm{dlog}_{g} : (\mathbb{Z}/q\mathbb{Z})^{\times} \;\xrightarrow{\ \sim\ }\;
  \mathbb{Z}/\varphi(q)\mathbb{Z},
  \qquad g^{\ell} \longmapsto \ell,
  \qquad \mathrm{dlog}_{g}(xy) = \mathrm{dlog}_{g}(x) + \mathrm{dlog}_{g}(y) ,
\end{equation}
and the characters of the unit group are
\begin{equation}
  \label{eq:character}
  \chi_{k}(x) \;=\; \exp\!\Bigl( \tfrac{2\pi i\, k\, \mathrm{dlog}_{g}(x)}{\varphi(q)} \Bigr),
  \qquad k \in \mathbb{Z}/\varphi(q)\mathbb{Z} .
\end{equation}
The frequency index $k$ lives in $\mathbb{Z}/\varphi(q)\mathbb{Z}$, not in
$\mathbb{Z}/n\mathbb{Z}$, and the two are routinely confused. Where the unit group is not
cyclic, the structure theorem gives $G \cong \mathbb{Z}_{o_{1}} \times \cdots \times
\mathbb{Z}_{o_{r}}$ with generators $g_{1}, \dots, g_{r}$; every unit is uniquely
$x = \prod_{j} g_{j}^{\,e_{j}}$, and the \emph{product characters} are
\begin{equation}
  \label{eq:prodchar}
  \chi_{\mathbf{k}}(x) \;=\;
  \exp\!\Bigl( 2\pi i \sum_{j=1}^{r} \tfrac{k_{j} e_{j}}{o_{j}} \Bigr),
  \qquad \mathbf{k} \in \widehat{G} \cong \prod_{j} \mathbb{Z}_{o_{j}},
  \qquad |\widehat{G}| = |G| .
\end{equation}
The logit grid is then a $2r$-dimensional array. In both cases $\chi_{\mathbf{k}}(xy) =
\chi_{\mathbf{k}}(x)\chi_{\mathbf{k}}(y)$, so multiplication becomes addition of exponent
tuples, which is what makes a clock expressible at all. For a matrix $M \in
\mathbb{R}^{G \times D}$ indexed by units, the character transform is
\begin{equation}
  \label{eq:transform}
  \widehat{M}_{\mathbf{k}} \;=\; \sum_{x \in G} \overline{\chi_{\mathbf{k}}(x)}\, M_{x}
  \;\in\; \mathbb{C}^{D} .
\end{equation}

A statistic reported in raw residue coordinates on a multiplication run is not a weaker
measurement, it is a different one: the same model's top eight logit components read $17.0\%$ ``$a+b$'' in raw coordinates
and $98.4\%$ in discrete-log coordinates.

\paragraph{Invariance to the generator.} The primitive root $g$ in \eqref{eq:dlog} is
arbitrary (our implementation returns the smallest), so nothing claimed about a circuit
may depend on it.

\begin{proposition}[Generator equivariance]
\label{prop:generator}
Let $\varphi = \varphi(q)$ and let $t$ satisfy $\gcd(t, \varphi) = 1$. Replacing $g$ by $g^{t}$
sends $\mathrm{dlog}_{g} \mapsto t^{-1} \mathrm{dlog}_{g} \bmod \varphi$ and therefore permutes
the character index by $k \mapsto t k \bmod \varphi$. Consequently
\begin{equation}
  \label{eq:equivariance}
  \mathrm{Gini},\;\; \mathrm{PR},\;\;
  \mathcal{L}(L),\;\; \mathcal{L}(\Pi_{K} L),\;\; \mathcal{L}(\Pi_{K^{c}} L),\;\; p
  \quad\text{are invariant,}
  \qquad
  K \;\longmapsto\; tK \bmod \varphi .
\end{equation}
\end{proposition}

\begin{proof}
$k \mapsto tk$ is a bijection of $\mathbb{Z}/\varphi\mathbb{Z}$, so it permutes the amplitude
multiset $\{A_{k}\}$; \eqref{eq:gini} and \eqref{eq:pr} are symmetric functions of that
multiset. It also maps the diagonal to itself, $(k,k) \mapsto (tk, tk)$, so $\Pi_{K}$ and
$\Pi_{K^{c}}$ in \eqref{eq:ablation} are conjugated by a permutation and their losses are
unchanged. The $p$ of \eqref{eq:equivariance} is the exact permutation $p$, a function of those
losses and of the uniform distribution over sets of fixed cardinality, which the permutation
preserves. Its Monte-Carlo estimate $\hat{p}$ of \eqref{eq:permp} is therefore invariant in
distribution, not draw by draw, because a fixed random seed selects different characters
under a different generator.
\end{proof}

Gini, PR and the three losses are verified numerically to $10^{-10}$ at $113$, $121$ and $125$,
and the losses per cyclic component at the non-cyclic moduli. For $\hat{p}$ the check is that
the $p < 0.01$ verdict does not change. A key-frequency value is therefore not a property of a model. Only
the set up to multiplication by a unit modulo $\varphi$ is, and this paper states key sets in
that sense. Writing ``key frequencies $\{11, 29, 33, 49\}$ at $n = 121$'' names a representative of an orbit, not
an invariant.

\subsection{Sparsity statistics, and the three conventions in the literature}
\label{sec:methods-sparsity}

Let $M$ be an embedding or activation matrix indexed by residues. For frequency $k$ we form
the sine and cosine projections $s_{k}, c_{k}$ and combine them into a single
\emph{amplitude}
\begin{equation}
  \label{eq:amplitude}
  A_{k} \;=\; \sqrt{\lVert s_{k} \rVert^{2} + \lVert c_{k} \rVert^{2}},
  \qquad k = 1, \dots, \lfloor \varphi/2 \rfloor,
\end{equation}
dropping the DC component $k = 0$. Two statistics are computed from $\{A_{k}\}$, and they are
computed on different powers of it.

Gini is computed on the amplitude. With $A_{(1)} \le \cdots \le A_{(M)}$ the sorted
amplitudes,
\begin{equation}
  \label{eq:gini}
  \mathrm{Gini}(A) \;=\;
  \frac{\sum_{i=1}^{M} (2i - M - 1)\, A_{(i)}}{M \sum_{i=1}^{M} A_{(i)}}
  \;\in\; \bigl[0,\; 1 - \tfrac{1}{M}\bigr],
\end{equation}
attaining $0$ on a flat spectrum and $1 - 1/M$ when a single frequency carries everything.
The upper bound depends on $M$, so a Gini is comparable only across spectra of the same
length. This is one reason the random-orthogonal control below must preserve that length.

The participation ratio is computed on the energy, with
\begin{equation}
  \label{eq:pr}
  p_{k} \;=\; \frac{A_{k}^{2}}{\sum_{j} A_{j}^{2}},
  \qquad
  \mathrm{PR}(A) \;=\; \frac{1}{\sum_{k} p_{k}^{2}}
  \;=\; \frac{\bigl(\sum_{k} A_{k}^{2}\bigr)^{2}}{\sum_{k} A_{k}^{4}}
  \;\in\; [1, M] ,
\end{equation}
its standard definition. $\mathrm{PR} = M$ for a flat spectrum and $\mathrm{PR} = 1$ when one
frequency carries everything. It is an effective count of active frequencies, which is why it
belongs on energy: $p_{k}$ must be a probability over frequencies.

The two choices are not interchangeable and neither is a matter of taste. On one of our
$n = 113$ checkpoints, Gini on energy reads $0.902$ where Gini on amplitude reads $0.543$.
Feeding amplitude into the participation ratio gives $12.5$ where the energy convention gives
$4.31$. Each error was live in this project's code for months, and in the second
case it produced a failed prediction against a published value before it was found.

Three incompatible conventions appear in this literature, and their values are not
comparable. Writing $u$ for the spectrum being summarised,
\begin{equation}
  \label{eq:conventions}
  \underbrace{\mathrm{PR}(u) = \Bigl(\textstyle\sum_{k} p_{k}^{2}\Bigr)^{-1} \in [1, M]}_{\text{ours, and \citet{nguyen2026dlogclock}}}
  \qquad
  \underbrace{\mathrm{IPR}_{r}(u) = \bigl( \lVert u \rVert_{2r} / \lVert u \rVert_{2} \bigr)^{2r} \in (0, 1]}_{\text{\citet{doshi2024grokking}, with } r = 2}
\end{equation}
move in \emph{opposite} directions: ours decreases with sparsity, theirs increases, and
$\mathrm{IPR}_{2}(u) = 1/\mathrm{PR}(u)$ exactly when both are computed on the same
normalised energy. \citet{nanda2023progress} use a third. A sparsity number quoted without its
convention cannot be checked, and we state ours at every use.

\paragraph{Never one basis alone.} A purely additive signal scores $0.53$--$0.63$ in the
multiplicative basis, so a multiplicative Gini reported on its own is uninterpretable. Every
sparsity measurement in this paper is reported as a triple (multiplicative, additive, and a
random-orthogonal control), together with the count of key frequencies.

\subsection{Controls}
\label{sec:methods-controls}

The random-orthogonal control must rotate the residue axis, not the feature axis.

\begin{proposition}[Rotating features is a no-op]
\label{prop:noop}
Let $M \in \mathbb{R}^{G \times D}$ and let $Q \in \mathbb{R}^{D \times D}$ be orthogonal.
Then the amplitude spectrum \eqref{eq:amplitude} of $MQ$ equals that of $M$ exactly.
\end{proposition}

\begin{proof}
The transform \eqref{eq:transform} acts on the residue index, so
$\widehat{MQ}_{\mathbf{k}} = \widehat{M}_{\mathbf{k}} Q$. Since $Q$ is orthogonal it preserves
the Euclidean norm of each row, so
$\lVert \widehat{MQ}_{\mathbf{k}} \rVert = \lVert \widehat{M}_{\mathbf{k}} \rVert$ for every
$\mathbf{k}$, and $A_{\mathbf{k}}$ in \eqref{eq:amplitude} is a function of that norm alone.
\end{proof}

A control built that way does not look broken. It returns the signal exactly, which reads as
a result worth reporting (``the clock is no better than a random basis''). That is the
failure mode Proposition~\ref{prop:noop} exists to prevent. The control we use instead acts on
the \emph{residue} index,
\begin{equation}
  \label{eq:randorth}
  M \;\longmapsto\; Q^{*} M, \qquad Q \in \mathrm{U}(|G|) \text{ Haar-random},
\end{equation}
which preserves the spectrum length $M$ (so the Gini bound in \eqref{eq:gini} is unchanged)
and destroys the character structure. Both directions are asserted in the analysis code's
self-check.

Two further controls appear throughout. A \emph{shuffle} control permutes cells within the
object being measured, so it preserves its size and marginal distribution and destroys its
structure. A \emph{matched-failure} control is built only from runs that failed to
generalise. The last of these is worth stating as a rule. A binary accuracy gate answers ``is
this a data point''. It is the wrong instrument for ``is this a control''. One of our runs
ends at test accuracy $0.9779$, a model that has learned the circuit. Scored as an ungrokked
negative control, it read $74.8\%$ neuron tuning and destroyed the contrast the control
exists to provide. We therefore classify every run into three states, grokked,
near-grok and failed, and build controls only from the failed ones. This biases against our own
findings, which is the direction that is rarely checked.

\paragraph{Classifying a run.}
% STE descriptive
A run's endpoint is a window, never its last logged row. Under weight decay $1.0$, post-grok
accuracy excursions are common. Across $11$ long runs we observe $19$ excursions below $0.90$
after grokking. Every one is exactly one $200$-step logging interval wide, and $18$ of them
recover. The nineteenth merely begins at the last step at which an observation exists.

For the logged test accuracies $\mathrm{acc}_{1}, \dots, \mathrm{acc}_{S}$ of a run, we
define
\begin{equation}
  \label{eq:window}
  \bar{a} \;=\; \mathrm{median}\bigl( \mathrm{acc}_{S-9}, \dots, \mathrm{acc}_{S} \bigr),
  \qquad
  \text{state} =
  \begin{cases}
    \text{grokked}   & \bar{a} > 0.99, \\
    \text{near-grok} & 0.90 \le \bar{a} \le 0.99, \\
    \text{failed}    & \bar{a} < 0.90,
  \end{cases}
\end{equation}
and test accuracy is read by column name rather than by position, because our two
implementations write different quantities at the same index.

There are three states, not two: a binary gate turns every near-miss into whichever side it
falls on. The run above ends at $\bar{a} = 0.976$, so it is near-grok and not a
control. Controls are built only from runs in the \emph{failed} state. The threshold $0.99$ is
the grokking criterion of every training script. The threshold $0.90$ is the floor of the
excursion count above, and it was fixed before any control arm was scored.
% STE end

\subsection{The ablation protocol}
\label{sec:methods-causal}

Sparsity says a representation is concentrated. It does not say the concentration is used.

Let $L \in \mathbb{R}^{G \times G \times n}$ be the logit tensor over a grid of units,
re-indexed into character coordinates by \eqref{eq:transform} on both input axes, so that
$\widehat{L}_{\mathbf{k}, \mathbf{k}'}$ is its component at the character pair
$(\mathbf{k}, \mathbf{k}')$. For a set $S \subseteq \widehat{G} \times \widehat{G}$ define the
spectral projection $\Pi_{S}$ by
\begin{equation}
  \label{eq:project}
  \widehat{\Pi_{S} L}_{\mathbf{k}, \mathbf{k}'} \;=\;
  \begin{cases}
    \widehat{L}_{\mathbf{k}, \mathbf{k}'} & (\mathbf{k}, \mathbf{k}') \in S, \\
    0 & \text{otherwise,}
  \end{cases}
\end{equation}
and write $\mathcal{L}(\cdot)$ for cross-entropy against the true targets. The three
quantities we report are then
\begin{equation}
  \label{eq:ablation}
  \underbrace{\mathcal{L}(L)}_{\text{baseline}}
  \qquad
  \underbrace{\mathcal{L}(\Pi_{K} L)}_{\text{restricted}}
  \qquad
  \underbrace{\mathcal{L}(\Pi_{K^{c}} L)}_{\text{excluded}}
\end{equation}
for a key set $K$, where $K^{c}$ is its complement. \emph{Restricted} keeps only $K$;
\emph{excluded} removes exactly $K$ and keeps everything else. We report all three as absolute
losses side by side, following \citet{chughtai2023toy}.

\paragraph{What this intervention acts on.} $\Pi_{S}$ acts on the
\emph{logit tensor} of a trained model. The forward pass is not re-run, no weight is ablated
and no activation is patched. This is \citet{nanda2023progress}'s restricted and excluded
loss protocol, applied in a coordinate they did not need, and what it licenses is a statement
about the output: the logits are carried by $K$ and are destroyed when exactly $K$ is removed.
It is not a statement about which internal component computes them.
\citet{chen2026multiplication} name ablations \emph{and} activation patching as the evidence
their account lacks; this protocol supplies the first. The title, the abstract and the heading
of Section~\ref{sec:causal} therefore name the ablation rather than assert causation, and
where the claim ledger of Appendix~\ref{app:ledger} uses the word \emph{causal} it means this
protocol and no other.

\paragraph{The one place we do re-run the forward pass.}
% STE descriptive
At $n = 113$ and $n = 121$ only, the same projection $\Pi_{S}$ of \eqref{eq:project} is
applied to the \emph{embedding} $W_{E}$ rather than to the logit tensor. The edit is made in
the multiplicative character basis. The network then runs forward unmodified
(Section~\ref{sec:causal-internal}). Softmax attention and the ReLU MLP run on top of the edited weights. So, unlike a mask on the output, nothing here is
algebraically forced. Non-unit rows and the read-out token row carry no character index and
pass through untouched.

The experiment was pre-registered before the models it uses were retrained. Its three
thresholds are inherited from criteria this paper fixes elsewhere for other purposes. They are
the $100\times$ of Section~\ref{sec:strata}, the $0.90$ generalisation floor, and
$p < 0.01$. So none was chosen after a number was seen.
% STE end

The two nonlinearities lie downstream of the edit, so the result is not a rearrangement of
the output, and the scoping of the logit-tensor protocol does not apply to it. It applies
everywhere else in this paper. We keep the two apart deliberately: two moduli with a weight
intervention do not retroactively upgrade the other twenty-one, measured a different way.

% STE descriptive
\begin{figure}[t]
  \centering
  \includegraphics[width=\textwidth]{../figures/F3_ablation_scheme.pdf}
  \caption[The two projections, drawn on the plane they act on.]{\textbf{The two projections, drawn on the plane they act on.} Each panel is the
    character plane of the logits at $n = 121$ seed $0$, in discrete-log coordinates, with
    DC at the centre. Shade is $\log_{10}$ of the magnitude
    $\lVert \widehat{L}_{\kappa_{a}, \kappa_{b}} \rVert$ summed over the output axis.
    \emph{Left:} the model as trained. The energy that matters lies on the diagonal
    $\Delta$, which is where \eqref{eq:diagonal} says a character circuit must put it.
    \emph{Centre:} $\Pi_{K}$ keeps the $4$ key characters and their conjugates and zeroes
    the other $12{,}091$ pairs.\par
    Rings mark the survivors, because nine kept pixels out of $110^{2}$ would otherwise read
    as nothing there. \emph{Right:} $\Pi_{K^{c}}$ removes exactly those nine and keeps the
    rest. The cross-entropy under each panel is printed beneath it: restricted is at
    or below baseline, and excluded is seven orders of magnitude worse. So the key characters
    are both sufficient and necessary on this run. The figure regenerates from
    \texttt{src/viz/plots.py::ablation\_scheme}. The spectrum, the key set and all three
    losses come from \texttt{src/analysis/ablation.py}, the same code
    Section~\ref{sec:causal} reports.}
  \label{fig:ablation-scheme}
\end{figure}
% STE end

\paragraph{The statistic is a permutation $p$, not a ratio to a control mean.} Drawing
$R_{1}, \dots, R_{B}$ uniformly at random from the character sets of cardinality $|K|$, with
$B = 10{,}000$ and $b$ the number of draws at least as damaging as the key set,
\begin{equation}
  \label{eq:permp}
  \hat{p} \;=\; \frac{1 + b}{B + 1},
  \qquad
  b \;=\; \sum_{b'=1}^{B}
  \mathbf{1}\Bigl[\, \mathcal{L}(\Pi_{R_{b'}^{c}} L) \;\ge\; \mathcal{L}(\Pi_{K^{c}} L) \,\Bigr],
\end{equation}
the fraction of equal-sized random sets at least as damaging as the key set.
The $(1+b)/(B+1)$ form \citep{phipson2010permutation} is deliberate, and the naive $b/B$ is a trap we fell into. The
ratio $b/B$ reads exactly $0$ whenever no draw is as damaging as the key set, and a Monte-Carlo
$p$ of $0$ claims a resolution the design does not have. An earlier draft of this paper
printed $p = 0.0000$ to four decimals from a $200$-draw null whose smallest expressible
value is $1/201 \approx 0.005$. Here the floor is $\hat{p} = 1/10{,}001 \approx 1.0 \times
10^{-4}$, and where $b = 0$ the text states it in words as well as in the estimate. It is
threshold-free and scale-free, and it is the statistic every dependence claim in this paper
rests on. We previously reported
$\text{excluded} / \text{mean}(\text{control})$ over $20$ draws and have retired it. The
control distribution is bimodal: roughly $73\%$ of random character removals leave the loss
untouched, because the model does not use those characters, and the remainder are
catastrophic. Its mean is therefore decided by how many catastrophic draws happen to land, and
at $n = 121$ seed $0$ that ratio ranges over $21\times$ to $440\times$ across control random
seeds alone, on identical data. Where a descriptive ratio is useful we quote the median
control. The same permutation logic, at $20{,}000$ draws, is used for the enrichment test in
Section~\ref{sec:crt-law}.

\paragraph{Multiplicity.} This paper's permutation tests form a
family of $320$: $157$ measurable stratum-measurements in the primary arm and $26$ in the
engine arm, plus $137$ single-modulus ablations across eight sweeps and two runs of a retired engine arm. They test one
hypothesis repeatedly rather than searching for a significant one among many, but the count
is what it is and the reader should not have to reconstruct it. At $B = 10{,}000$ the floor
$1.0 \times 10^{-4}$ clears a Benjamini--Hochberg correction at $\alpha = 0.01$ and a
Bonferroni correction at $\alpha = 0.05$ (per-test $1.6 \times 10^{-4}$). It does not
clear Bonferroni at $\alpha = 0.01$, which would require a per-test level of
$3.1 \times 10^{-5}$ and hence $B \ge 32{,}000$. That comparison fails, and it is stated here
rather than omitted. The draw count was chosen for cost: $B = 10{,}000$ is about ten hours of
local CPU across the family and $B = 32{,}000$ about thirty. This reason is recorded here so
it is not later mistaken for a statistical judgement.

\paragraph{Two guards against circularity in the ablated set.}
For the unit-stratum results the key set is recovered independently of the ablation, by
embedding norm, and the set recovered by ablation rank agrees with it. For the
per-stratum results of Section~\ref{sec:strata} no key-frequency detector is used at
all: the ablated set is the \emph{whole diagonal}, fixed in advance by
Theorem~\ref{thm:stratum}. Selecting a set by ablation rank and then permutation-testing that
same set would be circular, and the second design excludes it by construction.

\paragraph{Three further statistics.} The \emph{enrichment} of a predicted
frequency set $P$ against a spectrum $A$ compares its mean mass to the global mean,
\begin{equation}
  \label{eq:enrich}
  \mathrm{enr}(P) \;=\;
  \frac{\frac{1}{|P|}\sum_{k \in P} A_{k}}{\frac{1}{M}\sum_{k} A_{k}} ,
\end{equation}
and is scored by the same permutation logic as \eqref{eq:permp} over $20{,}000$ random sets of
cardinality $|P|$. The \emph{diagonal energy share} of a stratum measures how much of the
logit energy sits on $\Delta$ of \eqref{eq:diagonal}, against the share an unstructured tensor
would give,
\begin{equation}
  \label{eq:dshare}
  D \;=\; \frac{\sum_{(\mathbf{k},\mathbf{k}') \in \Delta} \lVert \widehat{L}_{\mathbf{k},\mathbf{k}'} \rVert^{2}}
               {\sum_{\mathbf{k},\mathbf{k}'} \lVert \widehat{L}_{\mathbf{k},\mathbf{k}'} \rVert^{2}} ,
  \qquad
  R \;=\; \frac{D}{1 / (|G_{m}| + 1)} ,
\end{equation}
so $R$ is the ratio to exact chance. Finally a neuron $h : G \times G \to \mathbb{R}$ is
\emph{tuned} to frequency $\mathbf{k}$ when a single frequency's eight-bin family explains
most of its DC-free two-dimensional spectral energy,
\begin{equation}
  \label{eq:tuned}
  \mathrm{tune}(h) \;=\; \max_{\mathbf{k} \neq 0}
  \frac{\sum_{(\mathbf{j},\mathbf{j}') \in F(\mathbf{k})} \lVert \widehat{h}_{\mathbf{j},\mathbf{j}'} \rVert^{2}}
       {\sum_{(\mathbf{j},\mathbf{j}') \neq 0} \lVert \widehat{h}_{\mathbf{j},\mathbf{j}'} \rVert^{2}}
  \;>\; 0.85 ,
\end{equation}
where $F(\mathbf{k})$ is the eight-bin family $\{(\pm\mathbf{k}, 0), (0, \pm\mathbf{k}),
(\pm\mathbf{k}, \pm\mathbf{k})\}$. The threshold $0.85$ and the eight-bin family are
\citet{nanda2023progress}'s; scoring a single bin instead reads $9.6\%$ where the correct
statistic reads $85.0\%$ on the same model.

\paragraph{The two-dimensional bin convention, which is load-bearing.} On a grid indexed so
that multiplication is addition, a function of the product concentrates on the diagonal and a
function of the ratio on the antidiagonal:
\begin{equation}
  \label{eq:bins}
  f(a + b) \;\longleftrightarrow\; (k, k),
  \qquad
  f(a - b) \;\longleftrightarrow\; (k, -k).
\end{equation}
Reported the other way round, \eqref{eq:bins} inverts every conclusion drawn from it. We
verified it by planting signals of known frequency and locating them, never by deriving it.
Three separate bin-convention errors in this project were each found that way and none by
reasoning.

\paragraph{A generation diagnostic.} A key set that does not \emph{generate}
the character group cannot implement an exact clock, since two units on which every key
character agrees cannot be separated by any weighting of those characters. We report this as a
per-run diagnostic, and every key set recovered here generates its full group.

\subsection{What the measurements are worth}
\label{sec:methods-validity}

\paragraph{Noise in a single-checkpoint Gini.} Measured on three moduli at three
seeds over the final $12{,}000$ steps, $\mathrm{Gini}_{\text{mult}}$ has within-run standard
deviation $0.02$ to $0.044$. At $n = 125$ seed $2$ it wanders over $0.498$--$0.621$, a range of
$0.122$, with no direction. This is oscillation about a flat mean and not a trend. The
statistic plateaus by roughly $16{,}000$ to $20{,}000$ steps and then jitters. Against the
$0.10$ threshold we use for ``the clock survives'', a single checkpoint is therefore about a
$2\sigma$ test, and only about $1.6\sigma$ at $n = 125$. Where the comparison matters we
average over the final $10{,}000$ steps and report the within-run spread. It follows that part
of the across-seed spread reported for a single-checkpoint measurement is estimator noise
rather than seed variance. The two were not separated in the earlier measurements, and they are
not re-attributed here.

\paragraph{Convergence at forty thousand steps.} For a
statistic $\sigma$ sampled along a run at steps $t$, write its \emph{drift} against the final
checkpoint $T$ as
\begin{equation}
  \label{eq:drift}
  \mathrm{drift}(\sigma) \;=\;
  \frac{1}{|W|}\sum_{t \in W} \frac{\bigl|\sigma(t) - \sigma(T)\bigr|}{|\sigma(T)|},
  \qquad W = \text{the final } 10{,}000 \text{ steps},
\end{equation}
and call a seed converged when $\mathrm{drift} < 0.05$. Requiring that in at least four of
five seeds, $n = 113$
converges ($2.3\%$ participation-ratio drift, $5/5$ seeds) and $n = 121$ converges ($2.8\%$,
$4/5$). Four moduli do not: $n = 125$ at $5.6\%$ ($2/5$ seeds), $n = 119$ at $5.9\%$ ($3/5$),
$n = 120$ at $7.7\%$ ($3/5$) and $n = 165$ at $4.8\%$ ($3/5$). Every participation ratio at
those four is measured at a non-stationary point. Rather than re-measure at a longer horizon,
we retire the participation ratio at composite moduli and do not use it for cross-modulus
comparison there.
That a representation statistic need not be monotone in training is not a peculiarity of our
setup. \citet{demoss2025complexity} track a compression-based complexity across grokking and
find it \emph{rises} through memorisation and only then falls. The value read at any single
checkpoint therefore depends on where in that arc the checkpoint sits. Our drift criterion is the
narrow, per-statistic version of that caution, and the moduli it fails on are exactly the ones
we decline to compare.

\paragraph{The fibre residual.} Consequence~2 of Theorem~\ref{thm:stratum} says the target is
constant on the fibres of \eqref{eq:fibration}. To test whether the \emph{logits} are too, we
average a block $B$ over its fibres, write $\bar{B}$ for the result lifted back to the block,
and report the share of block variance that averaging destroys:
\begin{equation}
  \label{eq:resid}
  \mathrm{resid}(B) \;=\;
  \frac{\bigl\lVert B - \bar{B} \bigr\rVert_{F}^{2}}
       {\bigl\lVert B - \mathrm{mean}(B) \bigr\rVert_{F}^{2}} \;\in\; [0, 1].
\end{equation}
Averaging always smooths, so a small residual on its own means nothing. We therefore compare
against $\mathrm{resid}$ computed on
\emph{permuted} fibres of the same sizes, which have no algebraic meaning. This is the
size-matched defence that three retracted claims of ours needed and did not have.

\paragraph{Effect sizes for a pre-registered two-cell comparison.} Where an experiment asks
what fraction of a known gap a variable explains, we fix the two endpoints in the
pre-registration and report
\begin{equation}
  \label{eq:effectsize}
  f \;=\; \frac{\sigma_{\text{observed}} - \sigma_{\text{baseline}}}
                {\sigma_{\text{target}} - \sigma_{\text{baseline}}},
\end{equation}
with the decision rule committed before the observation exists: held at $f \ge 0.70$, not
held at $f \le 0.30$, and \emph{partial}, reported as such, in between.
Because the endpoints are fixed in advance, $f$ can exceed $1$ or fall below $0$, and we
report the value rather than clipping it.

\paragraph{Training precision changes an absolute sparsity number in PyTorch.} The same
PyTorch code at $n = 113$ reads $\mathrm{Gini}_{\text{mult}} = 0.5791$ at float32 and $0.7728$
at float64, while our from-scratch engine reads $0.7671$ and $0.7536$ at the two dtypes and is
effectively unaffected. The effect is an interaction between implementation and precision, not
a main effect of precision. We tested and retracted the main-effect claim, and
Appendix~\ref{app:retractions} records that. Every sweep in this paper is float32, and we
name the regime beside every absolute sparsity number. Comparisons between moduli, and
comparisons within one arm, are unaffected, which is what almost every result here rests on.

\section{Validation}
\label{sec:calib}

Every result below depends on an analysis pipeline being right. This section is the evidence
that it is, kept short.

\subsection{Calibration against published results}

Table~\ref{tab:calib} sets our measurements beside the published values they reproduce.

% STE descriptive
\begin{table}[t]
  \centering
  \caption{\textbf{Our measurements beside the published values they reproduce, under the
    same protocol.} The first two blocks are PyTorch float32 runs, and the third block is
    the from-scratch engine. A value after $\pm$ is the standard deviation across seeds.
    The reference values come from each paper's text, except in the $n = 165$ block. There
    the reference column is their published checkpoint, scored by our pipeline.}
  \label{tab:calib}
  \begin{tabular}{llll}
    \toprule
    & quantity & this work & reference \\
    \midrule
    \multicolumn{4}{l}{\citet{nguyen2026dlogclock}, $n = 113$, units only, their protocol, five seeds} \\
    & $\mathrm{Gini}_{\text{mult}}$ & $0.547 \pm 0.010$ & $0.579$ \\
    & $\mathrm{PR}_{\text{mult}}$ & $4.67 \pm 0.44$ & $4.1$ \\
    & $\mathrm{Gini}_{\text{add}}$ & $0.017 \pm 0.002$ & $0.071$ \\
    & $\mathrm{PR}_{\text{add}}$ & $55.80 \pm 0.05$ & $52.7$ \\
    & key frequencies & $4.6 \pm 0.8$ & $4$ \\
    \midrule
    \multicolumn{4}{l}{\citet{chen2026multiplication}, $n = 165$, our model against their checkpoint} \\
    & $\mathrm{Gini}_{\text{mult}}$ & $0.629$ & $0.637$ \\
    & participation ratio & $4.76$ & $4.85$ \\
    & key set against $P(n)$ & exact & exact \\
    \midrule
    \multicolumn{4}{l}{\citet{nanda2023progress}, $a + b \bmod 113$, one seed} \\
    & grokking step & $6{,}900$ & $5$--$10$k \\
    & key frequencies in $W_{E}$ & $5$ & $5$ \\
    & loss without the non-key frequencies & improves $7.3\times$ & improves \\
    & single-frequency neurons of $512$ & $85.0\%$ & $84.6\%$ \\
    \bottomrule
  \end{tabular}
\end{table}
% STE end

\paragraph{The strongest item runs on someone else's weights.} We ran our analysis code on
\citet{chen2026multiplication}'s published $n = 165$ checkpoint and on a model we trained
ourselves, and they agree (Table~\ref{tab:calib}). Two independently trained models, two
codebases, the same numbers. Their own $\mathcal{J}$-class claim also replicates on their
checkpoint: the embedding separation ratio is $1.03$ against $0.63$ for a shuffled-label
control. It is a modest separation whose meaning comes entirely from the control.

\paragraph{We replicate the multiplicative-grokking result, with two qualifications.}
Table~\ref{tab:calib} gives the five-seed values at $n = 113$ under
\citet{nguyen2026dlogclock}'s exact protocol. The first qualification is that their
key-frequency count does not replicate as a property of the arm. We find $4.6 \pm 0.8$ key
frequencies ($4, 6, 4, 4, 5$), matching their $4$ in three of five seeds. An earlier version of
this claim reported ``$4$ key frequencies, matching their $4$'' from seed $0$. That is a seed
property, and we withdraw the clause. The second qualification is that the agreement holds
because both arms are float32. The same code at float64 reads $0.7728$. That is not a claim
their number is wrong: it reproduces. But a published float32 Gini is not comparable to a
float64 one.

\paragraph{The engine reproduces \citet{nanda2023progress} on addition.} As a gate before any
multiplication result was trusted, we trained $a + b \bmod 113$ on our from-scratch autograd
engine. It grokked at step $6{,}900$, inside their stated $5$--$10$k band at weight decay
$1.0$, with five key frequencies in $W_{E}$ against their five. Ablating all non-key
frequencies \emph{improves} loss $7.3\times$ where they report an improvement, and $85.0\%$ of
$512$ neurons are single-frequency tuned against their published $84.6\%$. That last number,
from an independent implementation with its own autograd, is the strongest single piece of
evidence that the engine is sound.

\paragraph{Two of that gate's criteria were edited after they failed.} The pre-registered
criteria passed unmodified. Two additional criteria beyond them did not, and they were
rewritten to implement what \citet{nanda2023progress} state. One had scored a single
two-dimensional bin where they score a degree-two polynomial spanning eight. The corrected
version landing at $85.0\%$ against a published $84.6\%$ is independent evidence the fix is
right. The ordering was nevertheless wrong, and every criterion for every experiment after
that gate was frozen in a committed pre-registration before the run.

\paragraph{The gate-1 circuit is a clean Clock \textup{(1 seed)}.} Its top eight logit
components are all $a+b$ terms, ranked $(\pm 33, \pm 33)$ at $100\%$, $(\pm 45, \pm 45)$ at
$31.9\%$, $(\pm 5, \pm 5)$ at $28.0\%$ and $(\pm 1, \pm 1)$ at $24.0\%$, and \emph{the rank
order by logit energy is identical to the rank order by causal ablation delta}. Among the top
sixteen components the energy share is $100\%$ $a+b$ and $0\%$ $a-b$.

Critical dataset size is real here, reproducing \citet{power2022grokking}. The sweep in
Section~\ref{sec:clock-grok} is the measurement, and it is what identifies $49$, $54$ and $63$
at a $0.30$ training fraction as dataset-size nulls rather than algebraic obstructions.

\subsection{Softmax collapse}
\label{sec:calib-collapse}

Under cross-entropy with weight decay $1.0$, our engine exhibits the no-gradient regime
\citet{prieto2025grokking} describe. At $n = 17$ with a $0.40$ training fraction, the global
gradient norm falls $6.2$ orders of magnitude, from $3.41 \times 10^{-1}$ to
$2.35 \times 10^{-7}$. Meanwhile the mean probability of the correct class rises to
$1 - 3.4 \times 10^{-8}$ and the maximum logit sits near $50$. Past that point only weight
decay moves the weights. In float64 the probability never reaches $1$ exactly. It displays as
$1.000000$ from step $1{,}000$ and does not saturate in the arithmetic, so the no-gradient
regime here is driven by the logit gap rather than by a rounded probability.

\paragraph{StableMax mitigates the mechanism.} Trained identically, the StableMax arm retains
more gradient ($7.51 \times 10^{-7}$ against $2.35 \times 10^{-7}$ at step $4{,}000$), and its
correct-class probability stops two orders further from $1$, at $0.999998$ against
$1 - 3.4 \times 10^{-8}$. An earlier internal measurement of ours put the gradient ratio at
roughly $10\times$ ``at the same logit scale''. Re-run with a fixed seed and a saved artifact,
it is $3.2\times$. The two arms are also \emph{not} at the same logit scale: $50.1$ against
$19{,}154.3$, a factor of $383$, because StableMax's logits have no saturation ceiling. The
direction and the mechanism hold, and the ratio was never a like-for-like comparison. Neither
arm generalises at this configuration, which is critical dataset size at $n = 17$ and not a
StableMax failure.

\subsection{Two implementations, one mechanism}
\label{sec:impl}

Every sweep in this paper was run in PyTorch. The architecture is also implemented on an
autograd engine written from scratch for this project, and the two are compared here because
a mechanism that appears in only one implementation is a property of that implementation.

\paragraph{The engine is correct.} Against PyTorch at identical weights it agrees to
$4.4 \times 10^{-16}$ on the forward pass and $2.68 \times 10^{-15}$ on gradients, verified by
a finite-difference suite over every operator. It groks at four moduli ($113$, $119$, $121$
and $125$) at three seeds, and so reproduces the qualitative result of
Sections~\ref{sec:clock} and \ref{sec:strata} independently of PyTorch's autograd.

\paragraph{The two disagree about the magnitude of sparsity, and the cause is training
precision acting on PyTorch.} At byte-identical hyperparameters the engine's circuits read
$\mathrm{Gini}_{\text{mult}}$ about $0.2$ higher than PyTorch's at every modulus, $7$ of $7$
matched grokked pairs, with a smaller $\lVert W_{E} \rVert$ as well. The $2 \times 2$ resolves
it (Figure~\ref{fig:precision}):

\begin{center}
\begin{tabular}{lcc}
  \toprule
  & float32 & float64 \\
  \midrule
  from-scratch engine & $0.7671$ & $0.7536$ \\
  PyTorch & $\mathbf{0.5791}$ & $\mathbf{0.7728}$ \\
  \bottomrule
\end{tabular}
\end{center}

Each cell is tail $\mathrm{Gini}_{\text{mult}}$ at $n = 113$: the mean over the final
$10{,}000$ steps, never a single checkpoint, over three seeds of $40{,}000$ steps. The experiment
was pre-registered before either dtype was read. Every cell is recomputed from its artifact
rather than quoted.

Precision is inert in the engine and decisive in PyTorch. Take $\sigma_{\text{baseline}}$ to be
PyTorch at float32 and $\sigma_{\text{target}}$ the engine's float64 value. The effect size
\eqref{eq:effectsize} is then $f = +1.110$ on Gini and $f = +0.968$ on
$\lVert W_{E} \rVert$, both above the pre-registered threshold of $0.70$. The Gini value
exceeds $1$, meaning float64 moves PyTorch slightly \emph{past} the engine's position, and it
is given unclipped. Float64 also moves PyTorch's neuron tuning from $0.938$ to $0.770$. That
is both axes at once, onto the engine's joint position, which is what a single convergence
axis looks like. It dissolves an earlier worry that the two might be running different
algorithms.

\paragraph{This is an interaction, not a main effect, and the distinction is a retraction.}
We previously claimed that training precision changes measured sparsity, full stop, and tested
it by varying dtype \emph{inside the engine alone}. The effect size \eqref{eq:effectsize} came
back $f = -0.07$, the sign wrong and not merely the magnitude small, and the claim was
retracted. It stays retracted. What is true is the narrower sentence: training precision
changes the measured sparsity \emph{in PyTorch}.

Consequently every absolute sparsity number this paper reports from PyTorch is a float32 measurement,
and we name the regime beside each one. Comparisons between moduli, and comparisons within a
single arm, are unaffected; almost everything this paper claims is one of those. The
replication of \citet{nguyen2026dlogclock} in Section~\ref{sec:calib} matched their published
value because both arms are float32, and that agreement is real.

Before precision was reached, we eliminated the following candidates:
\begin{itemize}
  \item the measurement protocol;
  \item the hyperparameters, read from the kernel source rather than assumed;
  \item the tail window;
  \item the optimiser (the engine's AdamW is algebraically PyTorch's);
  \item gradients ($2.68 \times 10^{-15}$);
  \item minibatch noise (both are full-batch);
  \item the initialisation distribution;
  \item the train/test split (byte-identical, the same $3{,}830$ pairs);
  \item the execution substrate (PyTorch on CPU reads $0.5791$ against $0.5665$ on a T4, so
    no Kaggle number carries a hardware artefact);
  \item the specific initialisation draw.
\end{itemize}

\paragraph{What remains open is the engine's dtype-robustness, and it has no surviving
candidate.} The gap is explained; why the engine reaches the float64 answer at either dtype is
not. Under a pre-registration committed before any dtype was read we tested and eliminated
three mechanisms. They are residual type promotion in intermediate computations ($0$ of $115$
forward intermediates promoted, with a positive control at $100\%$), matmul accumulation order,
and softmax formulation (ratios $1.000$--$1.172$, inside the pre-registered equivalence band).
The question stays open with no candidate rather than with a proposed one.

% STE descriptive
\begin{figure}[t]
  \centering
  \includegraphics[width=0.62\textwidth]{../figures/F7_precision.pdf}
  \caption{\textbf{Training precision is an interaction, not a main effect.} Each bar is tail
    $\mathrm{Gini}_{\text{mult}}$ at $n = 113$ over three seeds and $40{,}000$ steps. The
    from-scratch engine reads essentially the same value at either dtype, and PyTorch moves
    by $0.19$. The experiment was pre-registered before either dtype was read. The
    figure regenerates from \texttt{src/viz/plots.py::precision\_2x2}.}
  \label{fig:precision}
\end{figure}
% STE end

\section{Grokking, and the clock at prime powers}
\label{sec:clock}

Table~\ref{tab:glance} lists the results of this section and the next three, each with the
criterion it was scored against and its status in the claim ledger of
Appendix~\ref{app:ledger}.

% STE descriptive
\begin{table}[t]
  \centering
  \footnotesize
  \caption{\textbf{The results at a glance.} Each row gives a claim, the section that
    reports it, its criterion, the outcome and the scale of the evidence. The status column
    is copied from the claim ledger of Appendix~\ref{app:ledger}. A criterion marked
    ``pre-reg.'' was committed before the data it scores existed.}
  \label{tab:glance}
  \setlength{\tabcolsep}{4pt}
  \begin{tabular}{l p{0.24\textwidth} l p{0.19\textwidth} p{0.19\textwidth} l}
    \toprule
    & claim & \S & criterion & outcome & status \\
    \midrule
    C3 & models grok at most moduli & \ref{sec:clock-grok} & test acc.\ $> 0.99$ & $19$ of $23$ moduli & verified \\
    C6 & the clock survives prime powers & \ref{sec:clock-primepowers} & $\Delta < 0.10$, pre-reg. & five cyclic prime powers & verified \\
    C35 & $n = 113$ is not special & \ref{sec:clock-primes} & $\Delta < 0.10$, pre-reg. & $\Delta \le 0.0250$ at $127$, $131$ & verified \\
    C31 & the clock is in the neurons & \ref{sec:clock-neurons} & $\mathrm{tune} > 0.85$, pre-reg. & $K_{h} = K_{W}$ in $14$ of $14$ runs & verified \\
    C21 & the circuit uses the key frequencies & \ref{sec:clock-circuit} & top-eight $a+b$ energy & $98.37$--$100.00\%$ & verified \\
    C22 & logit ablation at prime powers & \ref{sec:causal-primepowers} & $\hat{p} < 0.01$ & $5/5$ at six moduli & verified, 5 seeds \\
    C23 & the same with no discrete log & \ref{sec:causal-nodlog} & $\hat{p} < 0.01$ & $28$ of $29$ runs & verified \\
    C36 & the same at $2^{7}$ & \ref{sec:causal-128} & $\hat{p} < 0.01$, pre-reg. & $\hat{p} = 1.0 \times 10^{-4}$ in $2/2$ & verified \\
    C32 & the same at three initialisations & \ref{sec:causal-init} & $\hat{p} < 0.01$ & $24/24$ & verified \\
    C39, C40 & key characters in the weights & \ref{sec:causal-internal} & three inherited criteria, pre-reg. & $3/3$ seeds at $113$ and $121$ & verified \\
    C33 & every stratum runs a clock & \ref{sec:strata-result} & G0 $\wedge$ G2 $\wedge$ G5, pre-reg. & $157$ measurements, $32$ strata & verified, 5 seeds \\
    C38 & the same clock across strata & \ref{sec:strata-sameness} & agreement $\ge 0.80$, pre-reg. & $158/160$ pairs & verified, 5 seeds$^{\dagger}$ \\
    C8 & additive key frequencies are $P(n)$ & \ref{sec:crt-law} & permutation $p < 0.01$, pre-reg. & $9$ of $9$ moduli, $3/3$ seeds & verified \\
    C37 & $\omega(n)$ at matched zero-divisor density & \ref{sec:crt-omega} & confirmatory, pre-reg. & $+0.804$ and $+0.578$ & confirmatory \\
    \bottomrule
  \end{tabular}
\end{table}
% STE end

\subsection{Where models grok, and why the exceptions are not algebraic}
\label{sec:clock-grok}

Under the single protocol of Section~\ref{sec:setup-training}, models grok in a majority of
seeds at $19$ of the $23$ moduli of Table~\ref{tab:moduli} (Figure~\ref{fig:grok}).
Generalisation to the held-out $70\%$ arrives long after the training split is fit. Only one modulus, $n = 49$, fails in
every seed. The four exceptions are $n = 49$ ($0/3$), $54$ ($1/3$), $63$ ($1/3$) and $91$
($1/3$).

Three of the four were pre-registered as dataset-size nulls before they were run, and a sweep
of the training fraction confirms it. At $\mathrm{train\_frac} \in \{0.30, 0.50, 0.80\}$:

\begin{center}
\begin{tabular}{lccc}
  \toprule
  $n$ & $0.30$ & $0.50$ & $0.80$ \\
  \midrule
  $49 = 7^{2}$ & 0/3 & 2/3 (+1 near-grok) & \textbf{3/3} \\
  $54 = 2 \cdot 3^{3}$ & 1/3 & \textbf{3/3} & \textbf{3/3} \\
  $63 = 3^{2} \cdot 7$ & 1/3 & \textbf{3/3} & \textbf{3/3} \\
  \bottomrule
\end{tabular}
\end{center}

The progression is monotone and no modulus resists. At $0.30$ these three hold $720$, $874$
and $1{,}190$ training examples, which is the critical-dataset-size regime
\citep{power2022grokking}; the obstruction is the size of the training split, not the algebra
of $n$. We state this because the alternative reading, that certain algebraic structures resist
grokking, is the interesting one, and it is wrong.

The fourth, $n = 91 = 7 \cdot 13$, is not explained by training-set size, and we record that
as a failed prediction of our own. Its modulus set was chosen by calibrating against the
observed transition at roughly $1{,}700$--$2{,}000$ training examples, and $91$ has $2{,}484$,
above $n = 81$, which groks $3/3$. It grokked $1/3$, with one further seed at
accuracy $0.976$, a near-grok that is counted as neither. Training-set size alone does not
predict grokability.

Grokking time is not interpretable from a single seed, and we do not interpret it. Across
five seeds at one modulus grokking time spans $6{,}600$ to $24{,}800$ steps, a factor of
$3.8$. One earlier claim of ours about the ordering of grokking times across moduli was
retracted for exactly this reason. A second was retracted because normalising grokking time by
$\varphi(n)$ manufactured an ordering. Raw steps \emph{decrease} with $\varphi$
($\rho = -0.665$), so dividing by $\varphi$ over-corrects and produces
$\rho(-\varphi, \text{steps}/\varphi) = +0.870$ over $15$ moduli. No algebraic property in
Table~\ref{tab:moduli} predicts grokking time in this data.

% STE descriptive
\begin{figure}[t]
  \centering
  \includegraphics[width=0.86\textwidth]{../figures/F8_grok.pdf}
  \caption[Generalisation, and the classification rule.]{\textbf{Generalisation, and the
    classification rule.} The panel plots test accuracy against step for the $30$ runs at
    the six moduli measured in Sections~\ref{sec:causal} and \ref{sec:strata}. Each run is
    coloured by the \emph{median of its final ten logged samples}, drawn as a ring at the
    right edge beside its final sample. Runs are classified into three states rather than
    two. A binary gate reports the one run at a window median of $0.976$ as a failure,
    where it is a near-grok. That model learned the circuit, and its use as an ungrokked
    control destroys the contrast a control exists to provide.\par
    The downward spikes are the $6$ post-grok excursions below $0.90$ that occur in $5$ of
    these runs. Each is exactly one $200$-step logging interval wide, and all six recover.
    None here is censored, but in the $120{,}000$-step horizon sweep one of $19$ such
    excursions began on a run's final logged sample. No later observation exists there, and
    an outcome read from that sample produced a claim we later retracted
    (Appendix~\ref{app:retractions}). The figure regenerates from
    \texttt{src/viz/plots.py::grok\_curves}.}
  \label{fig:grok}
\end{figure}
% STE end

\subsection{The discrete-log clock survives at prime powers}
\label{sec:clock-primepowers}

A prime power is the case \citet{chen2026multiplication} name: its ring has nilpotents and
non-regular $\mathcal{J}$-classes. If the character account required regularity, the
multiplicative-basis structure should degrade there. It does not.

On the unit sub-grid at five seeds, with the criterion fixed in advance: for a prime power
$n = p^{\alpha}$ with cyclic unit group,
\begin{equation}
  \label{eq:c6}
  \Delta(n) \;=\; \bigl| \mathrm{Gini}_{\text{mult}}(n) - \mathrm{Gini}_{\text{mult}}(113) \bigr|
  \;<\; 0.10 ,
\end{equation}
where each $\mathrm{Gini}_{\text{mult}}$ is the mean over seeds of the tail statistic. The
threshold is roughly $2$--$5$ times the estimator's own standard deviation
(Section~\ref{sec:methods-validity}), which is the only reason a single-checkpoint read would
have been too weak to use.

\begin{center}
\begin{tabular}{lccc}
  \toprule
  $n$ & $\mathrm{Gini}_{\text{mult}}$ & $\mathrm{Gini}_{\text{add}}$ & $|\Delta|$ vs $113$ \\
  \midrule
  $113$ (field) & $0.564 \pm 0.026$ & $0.018 \pm 0.004$ & --- \\
  $121 = 11^{2}$ & $0.556 \pm 0.028$ & $0.125 \pm 0.010$ & $\mathbf{0.008}$ \\
  $125 = 5^{3}$ & $0.558 \pm 0.048$ & $0.177 \pm 0.030$ & $\mathbf{0.005}$ \\
  \bottomrule
\end{tabular}
\end{center}

Mean $\pm$ standard deviation over five seeds; all values float32. The observed differences
are an order of magnitude below the threshold and well inside the seed spread. Three further
cyclic prime powers were added afterwards: $81 = 3^{4}$ at $3/3$ seeds within the threshold
($\Delta = 0.046, 0.024, 0.005$), $169 = 13^{2}$ at $2/3$ ($0.023, 0.013, 0.108$), and
$49 = 7^{2}$, which at $\mathrm{train\_frac}\ 0.50$ groks $3/3$ with
$\mathrm{Gini}_{\text{mult}} = 0.539, 0.574, 0.526$ ($\Delta = 0.020, 0.015, 0.033$). The
clock therefore survives at five cyclic prime powers ($7^{2}$, $11^{2}$, $5^{3}$, $3^{4}$ and
$13^{2}$), which span nilpotency depth $2$ to $4$.

Two qualifications travel with this. The first is that there are five, not six:
$128 = 2^{7}$ has a non-cyclic unit group and no discrete log, so $\mathrm{Gini}_{\text{mult}}$ is undefined rather than
unmeasured there; it is reached in Section~\ref{sec:causal-128} in product-character
coordinates. The second is that $169$ holds at $2$ of $3$ seeds, not $3$: the seed that missed is
reported, rather than a criterion that would have admitted it.

% STE descriptive
\begin{figure}[t]
  \centering
  \includegraphics[width=0.78\textwidth]{../figures/basis_comparison.pdf}
  \caption{\textbf{The same model in two bases.} The two panels give the spectral amplitude
    at $n = 113$ in the multiplicative (discrete-log) basis and in the additive basis.
    Sparsity is a property of the coordinate, not of the model. The density that makes
    multiplicative grokking look structureless in raw residue coordinates is a basis
    artefact \citep{nguyen2026dlogclock}. The figure regenerates from
    \texttt{src/viz/plots.py::basis\_comparison}.}
  \label{fig:basis}
\end{figure}
% STE end

What the prime powers show is that cyclicity of the unit group, not primality and not
regularity, is what the clock needs. $121$ and $125$ are local rings carrying nilpotents,
and in the multiplicative basis they are as sparse as the field.

\subsection{Three primes, not one}
\label{sec:clock-primes}

Every ``prime versus composite'' statement in the literature on this task, and in earlier
versions of this work, rested on $n = 113$. Two further primes were trained under the same
protocol:

\begin{center}
\begin{tabular}{lccccc}
  \toprule
  $n$ & $\mathrm{Gini}_{\text{mult}}$ & $|\Delta|$ vs $113$ & neurons tuned (dlog) & raw & shuffled \\
  \midrule
  $127$ & $0.5562$ & $0.0103$ & $0.908 \cdot 0.887 \cdot 0.998$ & $0.000$ & $0.000$ \\
  $131$ & $0.5915$ & $0.0250$ & $0.963 \cdot 0.975 \cdot 0.975$ & $0.000$ & $0.000$ \\
  \bottomrule
\end{tabular}
\end{center}

Both inside the same $0.10$ threshold, $3/3$ seeds each, float32. Nothing about $n = 113$ is
special; the prime case is a case, not an instance.

\subsection{The clock is in the neurons, not only the embedding}
\label{sec:clock-neurons}

Every multiplicative statistic above is computed on the embedding $W_{E}$, which leaves open
the objection that the structure is present in the representation and absent from the
computation. It is not. Pre-registered before the activations were scored, we score each of
the $512$ MLP neurons by $\mathrm{tune}(h)$ of \eqref{eq:tuned}. That is the share of its
DC-free two-dimensional spectral energy, over the unit grid reordered by discrete log, that a
\emph{single} frequency's eight-bin family explains. A neuron is tuned when that share exceeds
$0.85$:

\begin{center}
\begin{tabular}{lll}
  \toprule
  $n$ & tuned neurons, per seed & \\
  \midrule
  $113$ & $94.9 \cdot 91.8 \cdot 100.0 \cdot 99.8 \cdot 98.0\%$ & (published reference $96.9\%$) \\
  $121 = 11^{2}$ & $74.0 \cdot 79.9 \cdot 89.3 \cdot 73.2 \cdot 78.9\%$ & \\
  $125 = 5^{3}$ & $76.0 \cdot 72.9 \cdot 71.7 \cdot 85.7\%$ & \\
  \bottomrule
\end{tabular}
\end{center}

The controls are what make this a measurement. The same neurons in \emph{raw} residue
coordinates are $0.0\%$ tuned at every modulus and every seed; a per-neuron cell shuffle,
which preserves the grid's marginal distribution and destroys its structure, gives $0.00\%$ at
every grid side. A matched-failure control built only from runs that did not generalise reads
$6.8\%$ (our engine at $n = 125$, final accuracy $0.6082$) and $0.0\%$ (at $n = 49$ and
$n = 54$, accuracies $0.15$--$0.27$). And writing $K_{W}$ for the key set recovered from the embedding and $K_{h}$ for the set of
frequencies the tuned neurons select,
\begin{equation}
  \label{eq:keyagree}
  K_{h} \;=\; K_{W} \qquad \text{in } 14 \text{ of } 14 \text{ grokked runs,}
\end{equation}
element for element. The neurons and the embedding therefore run the same clock rather than
two coincidentally sparse ones. By Proposition~\ref{prop:generator} this is a statement
about two routes agreeing \emph{inside} a coordinate system, and it survives a change of
generator; the common set moves, the agreement does not.

The prime powers do not match the prime, and we do not explain the gap. $72$--$89\%$
against $92$--$100\%$ is a real and consistent shortfall, reproduced across seeds, and none of
our controls accounts for it. It is reported here and listed in
Section~\ref{sec:limits}.

\subsection{The circuit uses the embedding's key frequencies}
\label{sec:clock-circuit}

Finally, the frequencies identified in the embedding are the frequencies the output
uses. In discrete-log coordinates the top eight components of the logit two-dimensional DFT
are $a+b$ components. They carry $98.37$, $100.00$, $100.00$, $100.00$ and $100.00\%$ of
top-eight energy across five seeds, and they sit on the same frequency set the
$5\times$-median detector recovers from $W_{E}$. At seed $0$ the eighth component is an $a-b$ term carrying
$1.6\%$; we report the seed rather than the population value, because an earlier version of
this claim quoted $100\%$ from seed $0$, where the true figure is $98.37\%$.

The bin convention \eqref{eq:bins} is load-bearing here and was verified against planted
signals rather than derived; reported the other way round, the conclusion inverts.

\section{The clock carries the logits}
\label{sec:causal}

This is the section \citet{chen2026multiplication} name as missing. Their evidence for
stratified Fourier computation is, in their own words, \emph{``only correlational, meaning
causation cannot be directly inferred''}, and the future work they propose is exactly
\emph{``causal intervention techniques, such as ablations and activation patching''}. That is
why the ablation evidence here is a Results section rather than a line of support inside one.
\citet{nanda2023progress} and \citet{chughtai2023toy} fold ablations into their mechanism
sections, and we depart from that deliberately.

Sparsity says a representation is concentrated on a few characters. It does not say the
network computes with them. The protocol is Section~\ref{sec:methods-causal}'s: baseline, restricted
and excluded losses in exponent coordinates, against a null of $10{,}000$ random character
sets of the same cardinality, reported as a permutation $\hat{p}$.

\subsection{At prime powers}
\label{sec:causal-primepowers}

At five seeds across the six original moduli, not one of $10{,}000$ random equal-sized
character sets is as damaging as the key set in $26$ of the $31$ runs. The remaining five are
beaten by between $1$ and $2$ draws, so every run reads $\hat{p} \le 0.0003$:

\begin{center}
\begin{tabular}{llccc}
  \toprule
  $n$ & $(\mathbb{Z}/n\mathbb{Z})^{\times}$ & baseline/restricted, median & range over seeds & $p < 0.01$ in \\
  \midrule
  $113$ & $\mathbb{Z}_{112}$ & $2.34\times$ & $1.56$--$264$ & $5/5$ \\
  $121 = 11^{2}$ & $\mathbb{Z}_{110}$ & $1.85\times$ & $0.349$--$4.22$ & $5/5$ \\
  $125 = 5^{3}$ & $\mathbb{Z}_{100}$ & $1.98\times$ & $1.46$--$2.87$ & $5/5$ \\
  $119 = 7 \cdot 17$ & $\mathbb{Z}_{6} \times \mathbb{Z}_{16}$ & $0.67\times$ & $0.0146$--$1.96$ & $5/5$ \\
  $120$ & $\mathbb{Z}_{2}^{3} \times \mathbb{Z}_{4}$ & $1.25\times$ & $1.18$--$1.75$ & $5/5$ \\
  $165 = 3 \cdot 5 \cdot 11$ & $\mathbb{Z}_{2} \times \mathbb{Z}_{4} \times \mathbb{Z}_{10}$ & $1.42\times$ & $0.000709$--$1.63$ & $5/5$ \\
  \bottomrule
\end{tabular}
\end{center}

The improvement factor is reported as a median with its range across seeds, never as a mean.
Its distribution is heavy-tailed for the same reason the control distribution is. At $n = 113$
one seed of five reads $264\times$ where the other four read $1.56$ to $8.0$, and a mean over
those is $55.5 \pm 104.1$. That summary's standard deviation is twice its value, and one run
decides it. The claim rests on the permutation $p$, which is $5/5$
everywhere and is not a ratio at all.

Figure~\ref{fig:causal} gives the absolute losses and the null itself. At $n = 121$
seed $0$ the baseline loss is $4.96 \times 10^{-7}$, keeping only the key characters
\emph{improves} it to $2.10 \times 10^{-7}$, and removing exactly those characters raises it to
$1.29 \times 10^{1}$, seven orders of magnitude.
Keeping only the key characters improves the loss in $25$ of the $30$ five-seed runs, and in
$5/5$ seeds at $113$, $120$ and $125$.

The prime powers are the point. $121$ and $125$ are the local rings with nilpotents
and non-regular $\mathcal{J}$-classes, and their clock frequencies are load-bearing,
not merely a sparse description of the embedding.

\subsection{Where there is no discrete log}
\label{sec:causal-nodlog}

$119$, $120$ and $165$ have non-cyclic unit groups: no generator, no discrete log. They still
have an exponent coordinate, multiplication is still addition of exponent tuples, and a
character is a multi-index. Ablating there separates the key set from the median control by $10.0$ to
$4.6 \times 10^{7}\times$, so the
product-character circuit carries the logits too.

Across eleven further moduli at up to three seeds each, the permutation test passes at
$p < 0.01$ in $28$ of $29$ analysable runs, including $81 = 3^{4}$ at $3/3$. We name the
exception rather than the rate, and it is marginal rather than absent: $n = 98$ is $2/3$, with
one seed at $\hat{p} = 0.0151$ ($150$ of $10{,}000$ draws as damaging as the key set).

This count moved when the null was deepened, and the pre-registration said in advance that it
might. At $200$ draws it read $27$ of $29$. The second exception was $n = 49 = 7^{2}$ at
$p = 0.010$, two draws out of $200$, which landed exactly on the threshold.
\texttt{PREREGISTER\_r1\_permutation\_B.md} named both marginal tests before the re-run and
committed to reporting whatever they became, \emph{including} $26$ or $28$. At $10{,}000$
draws $n = 49$ reads $\hat{p} = 0.0050$ ($49$ draws) and passes. The other end of the range
did not move: $n = 98$'s seed read $0.015$ from $3$ of $200$ draws and reads $0.0151$ from
$150$ of $10{,}000$.

The pass at $n = 49$ is a caveat about the statistic, not a result about $n = 49$. That run
\emph{did not grok}: its held-out accuracy over the final ten samples is $0.891$, and its baseline loss, $1.3 \times 10^{-1}$, is four orders of magnitude
worse than any grokked run here. An earlier draft read its $p = 0.010$ as the test
``correctly'' reporting the absence of a circuit. It was not doing that, and the $50\times$
resolution is what exposed the reading. A permutation test on the logit spectrum asks whether
the key characters carry the logits; a model that has memorised much of its training grid
still has a logit tensor that is a function of $u \cdot v$ where it was trained. This is the
same limit Gate~2 measures directly. There $G1$, the same statistic and the named primary
criterion, passes on $20$ of $20$ \emph{failed} stratum-measurements
(Section~\ref{sec:strata-notclaimed}, Figure~\ref{fig:discrimination}). The permutation $p$
is evidence that a circuit is load-bearing \emph{where one exists}. It is not a detector of
whether the network generalised, and nothing in this paper asks it to be.

The analysis code was validated on a planted signal at $n = 35$ before it was trusted on
data.

\subsection{At $2^{7}$, the deepest nilpotency tested}
\label{sec:causal-128}

$n = 128$ is the sixth prime power and the hardest case in the set: $\omega = 1$, half its
residues zero divisors, unit group $\mathbb{Z}_{2} \times \mathbb{Z}_{32}$ with no discrete
log, and nilpotency depth $7$, the deepest tested. In product-character coordinates the key
characters carry the logits at the floor of the null, $\hat{p} = 1.0 \times 10^{-4}$, in $2/2$
analysable runs: $0$ of $10{,}000$ draws is as damaging as the key set in either.

The permutation $p$ is the claim here, and the improvement ratio is not. The ratio is given as
a median with its range. These samples are two or three runs drawn from a distribution
this paper has already shown to be heavy-tailed (Section~\ref{sec:methods-causal}), and a mean
over two or three such draws is exactly the summary we retired. Writing
baseline/restricted per run:

\begin{center}
\begin{tabular}{lcccc}
  \toprule
  $n$ & runs & median & range & $p < 0.01$ \\
  \midrule
  $128 = 2^{7}$        & $2$ & $239.52\times$ & $2.77$--$476.26$ & $2/2$ \\
  $131$                & $3$ & $2.88\times$   & $2.44$--$781.26$ & $3/3$ \\
  $127$                & $3$ & $5.04\times$   & $2.18$--$51.63$  & $3/3$ \\
  $91 = 7 \cdot 13$    & $2$ & $117.26\times$ & $1.45$--$233.07$ & $2/2$ \\
  $123 = 3 \cdot 41$   & $3$ & $1.76\times$   & $0.00$--$2.19$   & $3/3$ \\
  \bottomrule
\end{tabular}
\end{center}

\noindent
$13$ of $13$ runs across these five moduli reach the floor, $\hat{p} = 1/10{,}001 = 1.0 \times
10^{-4}$: no draw of $10{,}000$ was as damaging as the key set in any of them.
The ranges are the finding, and they are wide: at $n = 131$ one run of three improves by $781\times$ and the
other two by under $3\times$, so a mean of $262\times$ describes no run in the sample. An
earlier draft of this section reported exactly those means with standard deviations of the
same order as the values, which is the defect Section~\ref{sec:methods-causal} retires two
sections earlier. The ratios above supersede them. What survives unchanged is the
threshold-free statistic: every one of the $13$ runs is beyond every one of its $10{,}000$
draws.

Two qualifications. $n = 128$ grokked $2$ of $3$ seeds, so the $2/2$ is two analysable runs and
not two of three. And $n = 91$'s count is not counted as independent support. Our two analysis paths disagree
on its seed count, because one admitted a run at accuracy $0.976$. Under the three-state rule
that run is a near-grok and is neither data nor control.

The prime-power set therefore reaches $p = 2$ and depth $7$. The discrete-log
statistics of Section~\ref{sec:clock-primepowers} cover five cyclic prime powers; the ablation
covers all six.

\subsection{Independent of initialisation}
\label{sec:causal-init}

The dependence is not an artefact of one initialisation scale. Across $24$ runs at $n = 113$,
in three initialisation conditions (balanced fan-in, upstream-heavy and downstream-heavy,
following \citet{kunin2024getrich}) at eight seeds each, the permutation test passes at
$p < 0.01$ in $\mathbf{24/24}$. Every one of the $24$ is beyond \emph{every} one of its
$10{,}000$ draws.

\subsection{Inside the network, at a prime and at a prime square}
\label{sec:causal-internal}

Everything above acts on the logit tensor. It establishes what the output depends on, and
Section~\ref{sec:methods} states that it is not a statement about which internal
component computes it. At $n = 113$ and at $n = 121 = 11^{2}$ we can say the stronger thing.
The protocol, and the three criteria it inherits, are stated in
Section~\ref{sec:methods-causal}.

% STE descriptive
\begin{table}[t]
  \centering
  \small
  \begin{tabular}{cccccrl}
    \toprule
    $n$ & seed & baseline & restricted (acc.) & excluded (acc.) & excl./base. & $\hat{p}$ \\
    \midrule
    $113$ & $0$ & $1.2102\mathrm{e}{-05}$ & $2.0775\mathrm{e}{-06}$ ($1.0000$) & $4.7136\mathrm{e}{+01}$ ($0.0089$) & $3.895\mathrm{e}{+06}$ & $1/10001$ \\
    & $1$ & $1.2743\mathrm{e}{-07}$ & $3.2178\mathrm{e}{-07}$ ($1.0000$) & $3.5425\mathrm{e}{+01}$ ($0.0089$) & $2.780\mathrm{e}{+08}$ & $1/10001$ \\
    & $2$ & $1.5480\mathrm{e}{-07}$ & $2.2235\mathrm{e}{-07}$ ($1.0000$) & $2.3883\mathrm{e}{+01}$ ($0.0091$) & $1.543\mathrm{e}{+08}$ & $5/10001$ \\
    \midrule
    $121$ & $0$ & $1.7482\mathrm{e}{-06}$ & $3.4061\mathrm{e}{-07}$ ($1.0000$) & $5.0365\mathrm{e}{+01}$ ($0.0000$) & $2.881\mathrm{e}{+07}$ & $1/10001$ \\
    & $1$ & $7.7458\mathrm{e}{-03}$ & $1.0560\mathrm{e}{-01}$ ($0.9596$) & $5.7783\mathrm{e}{+01}$ ($0.0120$) & $7.460\mathrm{e}{+03}$ & $1/10001$ \\
    & $2$ & $1.6364\mathrm{e}{-03}$ & $2.2674\mathrm{e}{-03}$ ($1.0000$) & $1.8467\mathrm{e}{+01}$ ($0.0000$) & $1.128\mathrm{e}{+04}$ & $3/10001$ \\
    \bottomrule
  \end{tabular}
  \caption[Deletion of the key characters from the embedding returns the network to
    chance.]{\textbf{Deletion of the key characters from the embedding returns the network to
    chance.} Each row gives cross-entropy on the unit grid after an edit of $W_{E}$ and a new
    forward pass, three seeds per modulus. At $n = 113$ excluded accuracy $0.0089$--$0.0091$
    \emph{is} chance, $1/112 = 0.00893$. At $n = 121$ it is $0.0000$, $0.0120$ and $0.0000$
    against chance $1/110 = 0.00909$ on the $110 \times 110$ unit grid. Restricted accuracy is
    $1.0000$ on five of six runs and $0.9596$ on the sixth. The excluded/baseline median is
    $1.543 \times 10^{8}$ (range $3.895 \times 10^{6}$--$2.780 \times 10^{8}$) at
    $n = 113$.\par
    At $n = 121$ it is $1.128 \times 10^{4}$ (range $7.460 \times 10^{3}$--$2.881 \times
    10^{7}$). Both are medians and ranges over three seeds, because a standard deviation over
    three points is not a statistic. The two moduli are never pooled. $\hat{p}$ is
    Phipson--Smyth $(1+b)/(B+1)$ \citep{phipson2010permutation} against $B = 10{,}000$ random character sets of the same
    cardinality. Seeds $0$ and $1$ sit at the design floor and are printed as the fraction.
    Four decimals cannot distinguish ``at the floor'' from ``just above it''.}
  \label{tab:internal}
\end{table}
% STE end

At $n = 113$ all three criteria hold on all three seeds, against a pre-registered requirement
of two of three. The checkpoints are not new models. They reproduce the archived runs bit-identically,
$\max\lvert \Delta W_{E} \rvert = \max\lvert \Delta W_{U} \rvert
= \max\lvert \Delta \mathrm{logits} \rvert = 0$, with grokking steps $4{,}500$, $6{,}500$ and
$9{,}700$ recovered exactly. So Table~\ref{tab:internal} describes the same networks the rest
of this section measures. The projections were checked against the spectrum and not only
against the loss. Restricted preserves key-bin amplitude to the digit and zeroes the
complement to $10^{-13}$, and the key set holds $81.6\%$ of embedding amplitude in $8$ of $56$
bins at seed $0$.

\paragraph{The null's upper tail reaches the effect, and the permutation $p$ is what shows
it.} On seed $2$ the control \emph{maximum}, $2.4464 \times 10^{1}$, exceeds the excluded loss
$2.3883 \times 10^{1}$, and on seed $0$ it reaches $99.2\%$ of it. A ratio to the control
median would read $8.2 \times 10^{7}$ and hide this entirely. It is the same heavy-tail hazard
that retired the control-mean ratio, in a new arm. The cause is the design and it is
conservative: the draw pool includes the key characters, so a draw can partially reproduce
the intervention. Post hoc, and not part of the pre-registered analysis, the damage resolves into a
dose-response that is monotone \emph{in the mean} with how many key characters a draw removes.
Removals of $0$, $1$, $2$, $3$, $4$ and $5$ give $1.26 \times 10^{-5}$, $4.95$, $10.3$, $16.5$,
$23.9$ and $30.0$ at seed $0$. Figure~\ref{fig:internal} shows that the within-level spread is
wide, though, so the dose-response describes means and not individual draws. A draw containing \emph{no} key character does nothing at all, within $1.04$--$1.11
\times$ of baseline, and across the three seeds $12{,}468$ such draws never exceed
$1.05 \times 10^{-4}$ against an excluded loss of $24$--$47$. This is direct evidence for the
reading already given for the bimodality in Figure~\ref{fig:causal}: the harmless mode is the
draws the model does not use.

% STE descriptive
\begin{figure}[t]
  \centering
  \includegraphics[width=\textwidth]{../figures/F10_internal.pdf}
  \caption[The intervention acts on a weight, and the null's tail is draws that partially
    perform it.]{\textbf{The intervention acts on a weight, and the null's tail is draws that
    partially perform it.} \emph{Left:} baseline, restricted and excluded cross-entropy at
    $n = 113$ after an edit of $W_{E}$ and a new forward pass, three seeds. The panel uses the
    same three roles and hues as Figure~\ref{fig:causal}. Excluded is direct-labelled with its
    \emph{accuracy}, because that is the claim: $0.0089$--$0.0091$ is chance, $1/112$, not a
    degraded model. Restricted holds $1.0000$ on every seed.\par
    \emph{Right:} every one of the $30{,}000$ null draws, pooled over seeds and divided by its
    own run's baseline so the three share one axis. The draws are split by how many key
    characters each draw removed. Counts are above the axis, bars are medians, and the band is
    the real intervention, $3.895 \times 10^{6}$ to $2.780 \times 10^{8}$. The $12{,}468$ draws
    that remove none sit at the baseline line and never exceed $1.05 \times 10^{-4}$ in
    absolute terms. Damage rises monotonically in the mean with overlap, but the within-level
    spread is wide. A single key character can be worth four orders or almost nothing, so the
    dose-response is a statement about means, not about individual draws.\par
    This decomposition is post hoc and not part of the pre-registered analysis. The figure
    regenerates from \texttt{src/viz/plots.py::internal\_intervention}.}
  \label{fig:internal}
\end{figure}
% STE end

\paragraph{At a non-square-free modulus.} A prime is the weakest place to test this, because
every nonzero residue is a unit and the clause ``non-unit rows pass through untouched''
covers a single row. At $n = 121 = 11^{2}$ it covers eleven (zero and ten nilpotents),
and those rows form a non-regular $\mathcal{J}$-class at $d = 11$ (Section~\ref{sec:setup-strata}).
After the edit the network still holds eleven intact embedding rows and a whole non-regular
stratum. We ran the identical instrument, unmodified, under a separate pre-registration
committed before the retrain, with all three thresholds inherited from the $n = 113$ arm. The
pre-registration named $\text{excluded}/\text{baseline} < 100\times$ as the outcome that
would confine the internal claim to a prime. The worst of three seeds reads
$7.46 \times 10^{3}$, and all three criteria hold on all three seeds
(Table~\ref{tab:internal}). The checkpoints again reproduce the archived runs, with
$\max\lvert \Delta W_{E} \rvert = 0$ and grokking steps $10{,}700$, $10{,}300$ and $5{,}000$.

Three differences from $n = 113$ are reported, not interpreted. First, excluded accuracy falls
\emph{below} chance on two of three seeds ($0.0000$ against $1/110$) where at the prime it sat
on chance. No criterion reads accuracy against chance, so this is descriptive and would need
its own pre-registration to become a claim. Second, the null's tail again reaches the effect
on seed $2$. The control maximum $18.6887$ exceeds the excluded loss $18.4667$, and both
beating draws remove three of the four key characters. The $19{,}426$ draws that remove none
never exceed $1.71 \times 10^{-2}$ against an excluded loss of $18$--$58$. Third, the ratios
sit about two orders lower at the floor than at $n = 113$; the pre-registration declined to
predict that they would match, and we draw nothing from the comparison.

\paragraph{What this does and does not license.} At $n = 113$ and $n = 121$ the key
multiplicative characters are the representation the network \emph{computes with}, not merely
a basis its output is sparse in. That is the second of the two instruments
\citet{chen2026multiplication} name as missing, at two moduli and three seeds each, one of
them non-square-free with a non-regular stratum left intact. The logit-tensor result remains
the evidence at the other twenty-one moduli, and the scoping of Section~\ref{sec:methods}
stands there. Restricted is worse than baseline on two of three seeds at each modulus: by
$2.5\times$ and $1.4\times$ at $n = 113$, and by $13.6\times$ and $1.39\times$ at $n = 121$, at
accuracy $0.9596$ or above and at least two orders below excluded. This is the same
necessary-but-not-sufficient signature the non-cyclic case shows above, and a shortfall of the
key-set detector rather than of the circuit.

\subsection{Two caveats on scope}
\label{sec:causal-caveats}

Necessary is not the same as sufficient, and at two moduli we can only show the first. In the notation of \eqref{eq:ablation} the two are distinct claims:
\begin{equation}
  \label{eq:necsuf}
  \underbrace{\mathcal{L}(\Pi_{K^{c}} L) \;\gg\; \mathcal{L}(L)}_{K \text{ is \emph{necessary}}}
  \qquad\text{versus}\qquad
  \underbrace{\mathcal{L}(\Pi_{K} L) \;\lesssim\; \mathcal{L}(L)}_{K \text{ is \emph{sufficient}}} ,
\end{equation}
and only the first is what the permutation $p$ of \eqref{eq:permp} tests. At $n = 119$ the restricted loss is $2.0\times$ \emph{worse} than baseline: the six key
characters are necessary, since excluding them is catastrophic, but keeping only them does not
reconstruct the model. Both losses are around $10^{-5}$, six orders below the excluded loss,
so this is a shortfall in the frequency detector rather than a failed circuit, but it is not
the clean ``restricted improves'' signature the other runs show. $n = 123$ reproduces the same
pattern at a modulus that had never been seen ($p < 0.01$ in $3/3$, baseline/restricted median only
$1.76\times$, range $0.003$--$2.19\times$). The initialisation sweep shows it is condition-specific: $3/8$ under
upstream-heavy initialisation, $0/8$ under both of the others. It tracks the unbalanced
initialisation, not the modulus and not the group structure.

A frequency can sit in the embedding and do nothing. In our addition control at $n = 113$,
the frequency $k = 23$ has embedding norm $23.1$ against a noise floor near $1$, and a logit
ablation delta below $4 \times 10^{-10}$ (1 seed). It is present in
$W_{E}$ and absent from the circuit. \citet{nanda2023progress} report the same thing. The
consequence for method is that embedding-level frequencies are not circuit-level frequencies,
and the causal set is the one to report.

% STE descriptive
\begin{figure}[t]
  \centering
  \includegraphics[width=\textwidth]{../figures/F4_causal.pdf}
  \caption[The causal test and the null it rests on.]{\textbf{The causal test and the null
    it rests on.} \emph{Left:} absolute baseline, restricted and excluded cross-entropy per
    modulus, five seeds each. Bars are the median, and every seed is a dot. Removal of exactly
    the key characters costs three to seven orders of magnitude. A model with only them is at
    or below baseline. \emph{Right:} the permutation null for $n = 121$ seed $0$, $10{,}000$
    random character sets of equal cardinality, with the median control and the real key set
    marked.\par
    The null is bimodal. About three quarters of random removals leave the loss at baseline,
    because the model does not use those characters. The remainder cause damage but reach at most
    $10.9$. The key set is at $12.9$, beyond every one of the $10{,}000$ draws, so $\hat{p}$ is
    the floor, $1.0 \times 10^{-4}$. The bimodality is why a ratio to the control \emph{mean}
    is not a statistic here and was retired. The figure regenerates from
    \texttt{src/viz/plots.py::causal\_test}.}
  \label{fig:causal}
\end{figure}
% STE end

\section{Every stratum runs a clock, on its own local group}
\label{sec:strata}

Sections~\ref{sec:clock} and \ref{sec:causal} measure the unit sub-grid. At a field that is
almost the whole table; at $n = 120$ it is $7.1\%$ of it. This section measures the rest, and
it is the paper's central result.

\subsection{What the algebra says to measure}
\label{sec:strata-algebra}

Theorem~\ref{thm:stratum} determines the measurement completely, which is what makes it a
test rather than a search. Fix a stratum $J_{d} \times J_{e}$ and let $m, w, q, G_{m}$ be as in
\eqref{eq:mw} and \eqref{eq:localgroup}. By \eqref{eq:fibration} the target factors through the
fibration, so the measurement is fixed in three steps:
\begin{equation}
  \label{eq:pipeline}
  \underbrace{L\big|_{J_{d} \times J_{e}}}_{\text{block of logits}}
  \;\xrightarrow{\ \iota_{d}^{-1} \times \iota_{e}^{-1}\ }\;
  \underbrace{\widetilde{L}(u, v)}_{\text{unit coordinates}}
  \;\xrightarrow{\ \pi_{d} \times \pi_{e}\ }\;
  \underbrace{\overline{L}(\bar{u}, \bar{v})}_{\text{collapsed onto } G_{m} \times G_{m}}
  \;\xrightarrow{\ \eqref{eq:transform}\ }\;
  \widehat{\overline{L}}_{\kappa, \kappa'} ,
\end{equation}
after which we ask whether the energy sits on $\Delta$ of \eqref{eq:diagonal}. Each arrow is
forced: $\iota$ by Proposition~\ref{prop:torsor}, $\pi$ by \eqref{eq:fibration}, and $\Delta$
by consequence~3 of Theorem~\ref{thm:stratum}.

No key-frequency detector is used anywhere in this section. The ablated set in
\eqref{eq:ablation} is $K = \Delta$, the \emph{whole diagonal}, fixed in advance by
\eqref{eq:diagonal} with $|\Delta| = |G_{m}|$. The obvious objection to any ablation result is
circularity: choosing a set by ablation rank and then asking a permutation test whether that
set is unusually damaging. Here the set cannot be
chosen, because the algebra names it.

The fibre-constancy of the target is asserted at every stratum of all six moduli by the
analysis code rather than assumed.

\subsection{Result}
\label{sec:strata-result}

Criteria were committed before any per-stratum number existed. We measure $32$ strata across
six moduli, $157$ stratum-run measurements in total. They span $12$ distinct local groups
($\mathbb{Z}_{10}$, $\mathbb{Z}_{16}$, $\mathbb{Z}_{20}$, $\mathbb{Z}_{100}$,
$\mathbb{Z}_{110}$, $\mathbb{Z}_{112}$, $\mathbb{Z}_{2} \times \mathbb{Z}_{10}$,
$\mathbb{Z}_{4} \times \mathbb{Z}_{10}$, $\mathbb{Z}_{2}^{2} \times \mathbb{Z}_{4}$,
$\mathbb{Z}_{2} \times \mathbb{Z}_{4} \times \mathbb{Z}_{10}$,
$\mathbb{Z}_{2}^{3} \times \mathbb{Z}_{4}$ and $\mathbb{Z}_{6} \times \mathbb{Z}_{16}$),
and every one carries the same circuit (Figure~\ref{fig:strata}).

\paragraph{Which strata are measurable, and why the cut cannot launder a failure.} A stratum
is \emph{measurable} when two structural conditions fixed in the pre-registration both hold.
Its local group satisfies $|G_{m}| \ge 10$ (S-b: below that the eight-bin family cannot
discriminate), and the split leaves it at least $100$ held-out cells (S-c: fewer cannot
separate a clock from memorisation). Blocks with $d = n$ or $e = n$ are excluded outright,
since $J_{n} = \{0\}$ makes them identically zero. Both conditions are read off the
algebra of $n$ and off the train/test split, never off a model's output, so the cut is
settled before any accuracy is computed and an exclusion can never stand in for a failure.
For G0 in particular, were the cut a function of the outcome, G0 could not fail by
construction.

\begin{center}
\small
\begin{tabular}{l p{0.28\textwidth} l p{0.34\textwidth}}
  \toprule
  & criterion & threshold & observed \\
  \midrule
  G0 & held-out accuracy per block & $\ge 0.95$ & $1.0000 \pm 0.0000$ at every measurable stratum \\
  G1 & permutation $p$, diagonal vs $10{,}000$ sets & $p < 0.01$ in $\ge 4/5$ seeds & $0$ of $10{,}000$ draws in $133/157$; $\hat{p} \le 0.0023$ in all $157$ \\
  G2 & excluded loss / baseline & $\ge 100\times$ & $3.98 \times 10^{2}$ to $2.71 \times 10^{8}$ \\
  G3 & fibre residual \eqref{eq:resid} vs permuted control & $\le 0.5 \times$ control & real $0.000$--$0.132$, control $0.396$--$0.929$ \\
  G4 & diagonal energy share $D$ of \eqref{eq:dshare} & $R \ge 5.0$ & $D = 0.837$--$0.997$ \\
  G5 & tuned neurons $-$ per-stratum shuffle & $\ge 0.20$ & $85$ of $87$ scored; $0.193$--$1.000$ vs shuffle $0.0000$ \\
  G6 & stratum sets differ by the algebra of $n$ & chain vs lattice & see below \\
  \bottomrule
\end{tabular}
\end{center}

G5 is defined only where the tuning statistic is, which is $87$ of the $157$ measurable
stratum-runs (Section~\ref{sec:strata-notclaimed}). Two of those $87$ fall below its
threshold: $n = 165$'s $J_{5} \times J_{15}$ and $J_{15} \times J_{5}$, both at $0.1934$
against a $0.20$ bar. This is the lower end of the range reported above and is not a pass.

Measuring the non-unit strata raises the share of the multiplication table the mechanism
accounts for. Writing $\mathcal{M}$ for the set of measurable strata, that share is
\begin{equation}
  \label{eq:coverage}
  \mathrm{cov}(\mathcal{M}) \;=\;
  \frac{1}{n^{2}} \sum_{(d,e) \in \mathcal{M}} |J_{d}| \, |J_{e}|
  \;=\; \frac{1}{n^{2}} \sum_{(d,e) \in \mathcal{M}} \varphi(n/d)\,\varphi(n/e),
\end{equation}
using $|J_{d}| = \varphi(n/d)$ from \eqref{eq:jclass}. It rises from the unit stratum alone
($\mathcal{M} = \{(1,1)\}$, giving $\varphi(n)^{2}/n^{2}$) to:

\begin{center}
\begin{tabular}{lccl}
  \toprule
  $n$ & unit stratum alone & measured here & local groups \\
  \midrule
  $113$ & $98.2\%$ & $98.2\%$ & $\mathbb{Z}_{112}$ (there is no non-unit stratum) \\
  $121 = 11^{2}$ & $82.6\%$ & $\mathbf{97.7\%}$ & $\mathbb{Z}_{110} \rhd \mathbb{Z}_{10}$ \\
  $125 = 5^{3}$ & $64.0\%$ & $\mathbf{89.6\%}$ & $\mathbb{Z}_{100} \rhd \mathbb{Z}_{20}$ \\
  $119 = 7 \cdot 17$ & $65.1\%$ & $\mathbf{88.6\%}$ & $\mathbb{Z}_{6} \times \mathbb{Z}_{16}$, $\mathbb{Z}_{16}$ \\
  $120$ & $7.1\%$ & $\mathbf{23.1\%}$ & $\mathbb{Z}_{2}^{3} \times \mathbb{Z}_{4}$, $\mathbb{Z}_{2}^{2} \times \mathbb{Z}_{4}$ \\
  $165 = 3 \cdot 5 \cdot 11$ & $23.5\%$ & $\mathbf{82.3\%}$ & $\mathbb{Z}_{2} \times \mathbb{Z}_{4} \times \mathbb{Z}_{10}$, and three more \\
  \bottomrule
\end{tabular}
\end{center}

% STE descriptive
\begin{figure}[t]
  \centering
  \includegraphics[width=\textwidth]{../figures/gate2_strata.pdf}
  \caption[Every measurable stratum, and the criteria it meets.]{\textbf{Every measurable
    stratum, and the criteria it meets.} Each row is one stratum $J_{d} \times J_{e}$, and the
    rows are grouped by modulus. $113$ contributes a single unit stratum, $121$ and $125$ the
    nested $p$-adic chains, $119$ and $165$ divisor lattices, and $120$ the widest lattice.
    Each row carries the stratum's local group $G_{m} = (\mathbb{Z}/(n/m)\mathbb{Z})^{\times}$
    and its score on each criterion of the table above. There are twelve distinct groups
    across the $32$ strata, from $\mathbb{Z}_{10}$ to
    $\mathbb{Z}_{2} \times \mathbb{Z}_{4} \times \mathbb{Z}_{10}$. The point of the figure is
    that the local group changes from row to row while the pattern of met criteria does
    not.\par
    The algebra of $n$ fixes how many strata there are and what group each carries, not
    whether a clock runs on it. Strata excluded by the measurability cut are not drawn, and
    the cut is structural and is defined above. The figure regenerates from
    \texttt{src/viz/plots.py::gate2\_strata}.}
  \label{fig:strata}
\end{figure}
% STE end

G6 is the second half of the claim. The stratum sets differ by the algebra of $n$ in
the way the algebra predicts: $121$ and $125$ give \emph{nested chains}, the $p$-adic
filtration; $119$, $120$ and $165$ give \emph{non-chain divisor lattices}; and $113$ gives a
single stratum. The mechanism does not change. The stratification does.

\subsection{The pre-registered primary criterion does not discriminate}
\label{sec:strata-falsifier}

The pre-registration named its own falsifier: a failed run must not show a local clock. It
fired. It is scored over $183$ grokked and $20$ failed stratum-measurements. These are the
$157$ measurable stratum-runs of the PyTorch sweep above, plus a further $26$ from the
from-scratch engine arm of Section~\ref{sec:strata-engine}, a larger population than the $157$
used for the criteria table.

% STE descriptive
\begin{figure}[t]
  \centering
  \includegraphics[width=\textwidth]{../figures/F9_discrimination.pdf}
  \caption[The pre-registered primary criterion does not separate, and the three that
    do.]{\textbf{The pre-registered primary criterion does not separate, and the three that
    do.} Each row is one Gate 2 criterion, scored on every measurable stratum of the grokked
    runs (filled, $183$) and the failed runs (open, $20$). The count is printed at each mark,
    and the rule between them is the criterion's power to discriminate.
    Rows are ordered by that gap. G0, G2 and G5 separate at $1/20$, $1/20$ and $1/14$. One
    failed measurement passes each, and it is the same one, named in the text.\par
    G3 and G4 separate partly. G1, the permutation $p$ named primary in the pre-registration,
    passes on $20/20$ failed measurements. Its two marks coincide, so the gap is zero. The
    verdict rests on G0 $\wedge$ G2 $\wedge$ G5. Counts differ by row, because G3 and G5 are
    undefined on strata with no fibre to collapse and on runs without saved activations. A
    row is scored only where it is defined.\par
    The figure regenerates from \texttt{src/viz/plots.py::gate2\_discrimination}, scored by
    \texttt{scripts/gate2\_discriminate.py}.}
  \label{fig:discrimination}
\end{figure}
% STE end

Once looked at, the explanation is not subtle: $30\%$ of every block is training data, so a model
that has merely memorised its training cells still has logits that are a function of
$u \cdot v$ there. G1 tests whether the logits are a function of the product, and
memorisation satisfies that. It is a necessary condition, not evidence. It was named primary
because it is the statistic this project had used elsewhere. That is the kind of inherited
authority that should have been checked against a failed-run control before it was promoted,
and now is.

The verdict stands, and its weight rests on G0 $\wedge$ G2 $\wedge$ G5, each of which
passes on exactly one failed measurement and no more. We do not write that all six criteria
passed.

That one measurement is worth naming rather than rounding away. It is the unit
stratum of $n = 49$, seed $2$, of the extended sweep, and it is the only measurable stratum
that run has. Its window-median test accuracy is $0.8908$, which clears our FAILED threshold
of $0.90$ by $0.009$, and its final ten logged samples rise, with one small dip, from $0.842$ to
$0.963$. It is a run still learning when the step budget ended, not one that plateaued having
learned nothing. On that stratum it reaches $0.959$ held-out accuracy and an
excluded/baseline ratio of $302$, so it passes G0, G2, G4 and G5 exactly as a grokked
measurement would. We place it in the failed arm because that is where our three-state
endpoint rule puts it. A FAILED threshold $0.01$ \emph{lower}, at $0.89$, would reclassify it as a near-grok, remove it from the arm, and restore the
$0/19$ this paper previously reported. The threshold is not tuned to this measurement: the
$0.90$ floor is inherited from the excursion analysis that retired our own de-grokking claim
(Appendix~\ref{app:retractions}), fixed a week before this arm was scored, and we keep it. The claim is therefore that G0, G2 and
G5 separate the two arms on every measurement but one borderline run, not that they separate
completely.

This is a narrowing, not a rescue, and the direction is the whole argument. The
decision rule was fixed in the pre-registration before the data existed: the fourth branch is
declared if G1 $\wedge$ G2 hold at every measurable non-unit stratum at two or more moduli and
G6 shows the stratum sets differ. That rule was met on its own terms: G1 \emph{held},
at $99\%$ of grokked stratum-measurements. What fired is a \emph{separate} clause of the same
document, the falsifier requiring that a failed run not show a local clock, and it is a finding
about what G1 can discriminate rather than about whether the rule was satisfied. The move we
then make, resting the weight on G0 $\wedge$ G2 $\wedge$ G5, \emph{discards} the
criterion our own pre-registration named primary and leaves a strictly smaller evidence base
than the one we registered. A post-hoc rescue selects criteria to save a claim; this
deselects one, after it passed, because it turned out not to separate. The same direction runs
through Appendix~\ref{app:retractions}: all five withdrawn claims were positive findings of
our own, three of them died to a size confound we went looking for, and not one retraction
replaced a claim with a stronger one.

\subsection{The same clock, or one clock each}
\label{sec:strata-sameness}

Every criterion above is computed inside one stratum. That cannot separate \emph{one
mechanism, indexed per stratum} from \emph{a private sub-circuit per stratum}: both put a
diagonal clock in every block. The discriminating question is whether two strata that share
a local group select the \emph{same characters} of it. The instrument for it is already in
this paper: \eqref{eq:keyagree} compares two independently recovered key sets element
for element on the unit stratum. Here it is applied between strata, pre-registered before
any cross-stratum number existed.

Which pairs can carry the test is fixed by the algebra, before any model is opened.
Two strata are comparable when their local groups are identical. $\mathcal{J}$-blocks
$(d,e)$ and $(e,d)$ are transposes of one \emph{commutative} task, so their diagonal spectra
agree for a reason that has nothing to do with a shared circuit; they are reported
separately and are never support. What remains is $28$ pairs at $n = 165$, $2$ at $n = 120$
and $2$ at $n = 119$, and \emph{none} at $n = 121$ or $n = 125$, whose only same-group
pair is a transpose. $n = 165$'s $G_{11}$, eight strata on one local group, carries the test.
Across five seeds that is $160$ pairs.

\begin{center}
\begin{tabular}{llcc}
  \toprule
  & statistic & observed & pre-registered \\
  \midrule
  \textbf{agreement} & key sets \emph{exactly} equal & $\mathbf{158/160 = 0.9875}$ & $\ge 0.80$ \\
  \textbf{null} & Monte-Carlo, key sets redrawn at the observed sizes & $\mathbf{1.0 \times 10^{-4}}$ & $< 0.01$ \\
  \textbf{control} & cell-shuffle, sizes kept, structure destroyed & $\mathbf{0/160}$ & $< 0.30$ \\
  transpose & descriptive only, never support & $54/54$ & --- \\
  \bottomrule
\end{tabular}
\end{center}

\noindent
Mean Jaccard overlap is $0.9982$ and the median key set has $2$ characters. The two
disagreements are both in $n = 165$ seed $1$, which reads $26/28$; the other four seeds read
$28/28$. The null is the floor a $10{,}000$-draw design can express, and it is not a formality:
at $G_{11}$ the folded spectrum has $5$ bins, so two independent $2$-character sets coincide
by chance with probability $1/\binom{5}{2} = 0.10$. The shuffle control does not fail by
selecting \emph{different} characters: $0$ of $160$ shuffled blocks has any key set at all,
so the shuffle removes the clock rather than relabelling it.

What this does not have is a failed-run control. The design called for the same statistic on
models that never generalised. There are exactly \emph{two} scorable pairs of that kind in the
whole project, and both come from one run at test accuracy $0.7502$. That model learned most
of the task, and our three-state rule nonetheless buckets it with the $0.146$ memorisers. Both agree, so the control
reads $2/2$ against its own $0.50$ bar; its own null reads $p = 1.0$, because two pairs
cannot discriminate anything. Every genuinely failed run either has no comparable pair by
algebra ($n = 49$, $63$, $125$) or produces \emph{no key set at all} ($n = 54$ seed $1$, at
accuracy $0.2733$). The control is underpowered rather than passing, and this section's claim
rests on the null and the shuffle control instead. The same shape, a statistic that is load-bearing elsewhere passing on failures here, is what
Section~\ref{sec:strata-falsifier} reports for the pre-registered primary criterion. It is
why we score every criterion on a failure when a failure exists.

\subsection{The single-seed precursor}
\label{sec:strata-c7}

The torsor coordinate was first sighted at one seed, before the pre-registration existed, by
indexing a non-regular $\mathcal{J}$-class not by a group law it lacks but by the unit action,
and taking the product-character transform in that coordinate. Every class beat its
random-basis control, regular or not: $J_{11}$ at $n = 121$ scored $0.345$ against $0.124$
($2.79\times$), $J_{5}$ at $n = 125$ scored $0.526$ against $0.156$ ($3.37\times$), and
$J_{2}$ at $n = 120$ scored $0.476$ against $0.162$ ($2.93\times$) (1 seed). That
observation is the origin of this section, and the section supersedes it: the same statement,
measured at five seeds across $26$ non-unit strata under criteria committed in advance.

An aggregate drawn from that single-seed table, that regular classes carry more structure
than non-regular ones, was retracted. It was confounded by class size, $J_{1}$ being always
regular and always largest. Size-matched within one model at $n = 120$, the only place such
pairs exist, the two overlap. Appendix~\ref{app:retractions} carries it with the other
retractions.

\subsection{Independent replication}
\label{sec:strata-engine}

The same verdict holds on our from-scratch engine at three moduli ($119$, $121$, $125$), two
to three seeds each: every measurable non-unit stratum carries a local clock. The two
implementations disagree about sparsity \emph{magnitude} for the reason given in
Section~\ref{sec:impl}, and agree about the stratified clock. One miss is reported:
the $n = 119$ unit stratum fails G1 in one of two seeds.

\subsection{What is not claimed}
\label{sec:strata-notclaimed}

\paragraph{G1, G2 and G4 are partly entailed by correctness.} By \eqref{eq:fibration} any
\emph{correct} function on a stratum is a function of $\mathrm{dlog}(\bar{u}) +
\mathrm{dlog}(\bar{v})$, hence supported on $\Delta$. So a perfectly confident lookup table
already reads $D = 1.000$, $\hat{p}$ at the floor and fibre residual $0$, and the analysis code
plants exactly that and measures it. What these criteria establish is that the logit tensor
depends on $(u,v)$ essentially only through $u \cdot v \bmod n/m$, i.e.\ that the computation
has collapsed onto the local ring. A lookup-table fit of $r^{2} = 0.012$--$0.074$ says the
logits are far from a scaled one-hot, so the collapse is a real statement about the tensor.
But a smooth function of $u \cdot v$ would also read $D \approx 1$, and these criteria cannot
by themselves say the mechanism is a \emph{character} circuit.

Sameness across strata is measured in Section~\ref{sec:strata-sameness}, and the limit of
that measurement is the absence of a failed-run control, not the absence of the
measurement. Every other statistic in this section is computed \emph{inside} one stratum,
against that stratum's own local group. A model that routed each stratum to a private
sub-circuit, each correct on its own block, would produce all of them.

G5 is the criterion that says ``clock''. Single-frequency MLP neurons in local-group
coordinates are entailed by nothing, and the shuffle control reads $0.0000$ at every stratum.
The tuning statistic as pre-registered is cyclic-only, so G5 covers $121$, $125$, $119$'s
sub-strata and $165$'s $\mathbb{Z}_{10}$ strata. The non-cyclic strata are covered by a
product-character version reading $0.182$--$0.877$ against a shuffle control of $0.0000$,
which was added after the pre-registration, is exploratory, and was never scored as G5.

The model does not solve every stratum. On held-out cells, $15$ stratum-runs fall to
$0.379$--$0.833$. Every one is below the measurability cut, and none affects the verdict. An
earlier internal record of ours claimed every stratum at accuracy $1.0000$, but that was the
full grid including training cells, and it is corrected here. The degradation is also not
evidence about the algebra. Every failing stratum has $|G_{m}| \le 8$ \emph{and} at most $64$
cells, and those are structurally coupled. So it cannot separate ``the local group is too
small to carry a clock'' from ``too few training examples''.

$n = 113$ contributes nothing to the stratum claim. A field has one stratum. It is
the positive control and that is all.

\section{The additive basis carries the CRT stratification}
\label{sec:crt}

The multiplicative basis is where the clock lives. The additive basis is not empty, and what
it holds is the stratification itself.

\subsection{The key additive frequencies are the CRT-component duals}
\label{sec:crt-law}

Let $n = \prod_{i=1}^{\omega(n)} q_{i}$ be the factorisation into \emph{maximal} prime powers
$q_{i} = p_{i}^{\,\alpha_{i}}$, so that the Chinese remainder theorem gives
$\mathbb{Z}/n\mathbb{Z} \cong \prod_{i} \mathbb{Z}/q_{i}\mathbb{Z}$. Define the predicted
additive frequency set
\begin{equation}
  \label{eq:crtset}
  P(n) \;=\; \bigcup_{i=1}^{\omega(n)}
  \Bigl\{\, t \cdot \tfrac{n}{q_{i}} \;:\; t \ge 1,\; t \cdot \tfrac{n}{q_{i}} \le \lfloor n/2 \rfloor \,\Bigr\}
  \;\subseteq\; \bigl\{1, \dots, \lfloor n/2 \rfloor \bigr\}.
\end{equation}
The ceiling at $\lfloor n/2 \rfloor$ is not cosmetic: the additive spectrum folds conjugate
frequencies, so it has exactly $\lfloor n/2 \rfloor$ bins and $k$ and $n - k$ are the same
frequency. Dropping the ceiling doubles $|P(n)|$ and the test no longer means anything. At
$n = 165$, \eqref{eq:crtset} gives $P = \{15, 30, 33, 45, 55, 60, 66, 75\}$; at $n = 120$ it
gives $\{15, 24, 30, 40, 45, 48, 60\}$.

\begin{quote}
The key additive frequencies of a grokked $a \cdot b \bmod n$ transformer are exactly $P(n)$:
one family per maximal prime power, the additive duals of the Chinese-remainder
components.
\end{quote}

The mechanism is elementary and needs no network. The indicator of the multiples of
$d$ has additive Fourier support on the multiples of $n/d$, and by \eqref{eq:jclass} each
$\mathcal{J}$-class is a union of arithmetic progressions generated by the CRT components.
At $\omega(n) = 1$ we have $n/q_{1} = 1$, so \eqref{eq:crtset} is every one of the
$\lfloor n/2 \rfloor$ bins: at $n = 113$, all $56$ of them. The prediction is then vacuous
and a prime power can never be counted as support.

Pre-registered before the moduli were run, and tested by a threshold-free permutation
enrichment over $20{,}000$ permutations. It holds at $9$ of $9$ testable moduli at $3$ of $3$
seeds each, every one at $p < 0.01$:

\begin{center}
\begin{tabular}{rcc}
  \toprule
  $n$ & enrichment, seeds 0/1/2 & seeds $p<0.01$ \\
  \midrule
  54 & 2.7/1.5/6.1 & 3/3 \\
  63 & 4.7/1.6/1.6 & 3/3 \\
  75 & 3.7/7.6/3.3 & 3/3 \\
  98 & 7.6/6.8/8.3 & 3/3 \\
  99 & 5.9/4.2/7.3 & 3/3 \\
  100 & 3.7/6.9/5.5 & 3/3 \\
  105 & 11.8/14.4/11.0 & 3/3 \\
  143 & 10.6/6.3/3.4 & 3/3 \\
  147 & 8.6/5.1/4.1 & 3/3 \\
  \bottomrule
\end{tabular}
\end{center}

% STE descriptive
\begin{figure}[t]
  \centering
  \includegraphics[width=\textwidth]{../figures/F6_crt.pdf}
  \caption[The predicted set is fixed by the factorisation before any model is trained, and
    the energy is on it.]{\textbf{The predicted set is fixed by the factorisation before any
    model is trained, and the energy is on it.} \emph{Top:} additive spectral energy of
    $W_{E}$ against frequency $k$, in raw residue coordinates, where $P(n)$ of
    \eqref{eq:crtset} lives. Bars in $P(n)$ are red, and the rest are grey. \emph{Bottom:}
    the threshold-free test. The grey histogram is the enrichment of $20{,}000$ random
    frequency sets of the same cardinality, and the red rule is the observed enrichment. The
    axis is logarithmic because the null concentrates near $1\times$.\par
    The rightmost column is \citet{chen2026multiplication}'s published checkpoint scored by
    our pipeline, which is the calibration for the other three. At $n = 119$ the effect is
    $3.30\times$ rather than $14.08\times$ and still lies beyond every draw. That modulus was
    enriched all along and invisible to a $5\times$-median threshold, which is why the test is
    threshold-free. Prime powers are omitted because $n/q_{1} = 1$ makes $P(n)$ every bin.
    That is vacuous by construction, not an absent measurement. The figure regenerates from
    \texttt{src/viz/plots.py::crt\_law}.\par
    The predicted set, the energy and the null come from \texttt{test\_crt\_law.py}. The
    generator asserts the published $14.08\times$ at $n = 165$ before it draws, because a
    bin-convention error here is silent at every modulus but one.}
  \label{fig:crt}
\end{figure}
% STE end

$8$ of those $12$ moduli (the nine of the table together with $165$, $120$ and $119$) are
non-square-free. Two further CRT-testable moduli were added
later: $n = 123 = 3 \cdot 41$ at enrichment $7.59 / 3.28 / 7.43$, $p < 5 \times 10^{-5}$ in $3/3$. At the
six original moduli the law holds at five seeds. $165$ reads $17.4, 13.3, 11.2, 11.2, 11.8$
and $120$ reads $13.8, 11.4, 9.8, 7.6, 13.3$, both $5/5$. $119$ reads
$3.3, 11.1, 5.2, 7.9, 6.4$, also $5/5$: it was enriched all along but invisible to a threshold
rather than absent.

The predicted and observed sets are \emph{exactly} equal, with no misses and no extras, at
$n = 165$ on both our own model and \citet{chen2026multiplication}'s published checkpoint, and
at $n = 120$. The second is the sharper test, because $120 = 2^{3} \cdot 3 \cdot 5$ has
generators $120/8$, $120/3$ and $120/5$: \emph{maximal prime powers, not primes}. Replacing
$q_{i}$ by $p_{i}$ in \eqref{eq:crtset} predicts $\{24, 40, 48, 60\}$ and misses
$\{15, 30, 45\}$, so the distinction is testable and the data decides it.

We asked the same question of the stratification with no network involved, and the answer
bounds the law rather than confirming it. The $\mathcal{J}$-class indicators of
$\mathbb{Z}/n\mathbb{Z}$ are sparse in the additive basis.
$\mathrm{Gini}_{\text{add}}$ is $0.406$ at $n = 121$, $0.450$ at $165$, $0.400$ at $120$, and
$\mathbf{0.000}$ at $n = 113$, a field having no stratification to encode. Their energy is
also strongly enriched on $P(n)$: $4.87\times$ at $n = 165$, $3.32\times$ at $120$ and
$7.15\times$ at $119$, every one beyond all $20{,}000$ permutations. But they do not
\emph{select} $P(n)$. Under the same detector and the same bin convention as the model
measurements, the indicators' key set equals $P(n)$ at only $2$ of the $10$ composite moduli
we tested. At $n = 165$ it holds $3$ of the $8$ predicted frequencies, and at $n = 120$ it is
$\{20, 40, 60\}$. Here $20 = 120/6$ is the dual of a divisor that is \emph{not} a maximal
prime power, so it lies outside $P(n)$ and outside the prime variant alike.

This is structural. The transform of $\mathbf{1}_{J_{d}}$ at frequency $k$ is the Ramanujan
sum $c_{q}(k) = \mu(q/g)\,\varphi(q)/\varphi(q/g)$, with $q = n/d$ and $g = \gcd(k, q)$. It
has magnitude at most $1$ wherever $k$ is coprime to $n$ and reaches $\varphi(q)$ where $k$
shares all of $q$. So the stratification is not confined to any divisor family: at
$n = 165$ its energy is nonzero at all $82$ frequencies. But it concentrates on the
frequencies that share a factor with $n$, the multiples of the duals of \emph{every}
divisor. At $n = 165$ those are $42$ frequencies, carrying $30.5\times$ the mean energy of
the other $40$, against $|P(165)| = 8$.

That makes the contrast with the network the result, and it runs in the paper's favour. On the same run, with one instrument: the model's key set is all $8$ of $P(165)$
and all $7$ of $P(120)$, with no misses and no extras, while the indicators select a subset
of the much larger divisor-dual family. Choosing the maximal-prime-power subfamily out of
the divisor duals is therefore something the trained network does, not something the algebra
hands it. The enrichment figures in this paragraph were computed after the set-equality test
they replace had failed, and are exploratory. The contrast table is not: both of its
columns use the detector fixed before the measurement.

The vacuous case excludes $113$, $121$, $125$, $127$, $131$, $128$, $49$, $81$ and $169$:
every prime and every prime power in our set.

\subsection{The early enrichment ramp is not the grokking circuit}
\label{sec:crt-ramp}

The CRT enrichment rises monotonically from early in training, it is specific to the CRT
structure against three size-matched nulls, and it would be a convenient progress measure. It
is not one, and this subsection is the negative result that says so.

We trained three seeds at $n = 63$ that did not generalise (test-accuracy window medians
$0.1801$, $0.1655$, $0.3183$, with training accuracy $1.000$ from step $2{,}000$, a clean
memorisation arm) and scored the same ramp, under criteria committed before the runs
existed:

\begin{center}
\begin{tabular}{lcccc}
  \toprule
  seed & final accuracy & enrichment at $1$k & enrichment at $20$k & ramp $\rho$ \\
  \midrule
  0 & $0.1803$ & $1.397$ & $\mathbf{1.653}$ & $\mathbf{1.000}$ \\
  1 & $0.1655$ & $1.339$ & $\mathbf{1.459}$ & $\mathbf{0.988}$ \\
  2 & $0.3177$ & $1.422$ & $\mathbf{1.655}$ & $\mathbf{0.995}$ \\
  \bottomrule
\end{tabular}
\end{center}

The failed runs show the same monotone ramp as grokked runs, and end \emph{above} the
$1.22$--$1.45$ band that grokked runs occupy early. Independent corroboration from PyTorch runs
at the same modulus reads $1.626$ and $1.574$ at accuracies $0.2236$ and $0.2749$. Two
implementations, two data sources, one answer.

The early ramp is therefore not evidence of circuit formation and must not be read as a
progress measure. This extends rather than contradicts the specificity result: that
already recorded that the permutation $p$ does not separate grokked from failed runs and that
only the enrichment \emph{magnitude} does. What is new is that the ramp's \emph{shape} does not
separate either, and monotonicity was the last property that might have made the early signal
diagnostic. Enrichment magnitude does still separate cleanly:
grokked runs reach $3.3$ to $17.4$ where these failures top out at $1.66$. The ablation
evidence for the CRT structure in grokked models (Section~\ref{sec:causal}) is untouched. Wherever the ramp appears in this
paper, the failed-run ramp appears beside it.

\subsection{Additive sparsity, zero divisors, and $\omega(n)$}
\label{sec:crt-omega}

Additive sparsity tracks zero-divisor density across $18$ moduli at $\rho = +0.858$. It also
separates into non-overlapping bands by $\omega(n)$, the number of distinct prime factors:
$\omega = 1$ spans $0.017$--$0.184$, $\omega = 2$ spans $0.394$--$0.481$, and $\omega = 3$
spans $0.578$--$0.589$, against a residue-axis random-orthogonal control flat at
$0.117$--$0.155$ throughout (Figure~\ref{fig:omega}).

This is a confirmatory result, not a discovery, and the distinction is not cosmetic.
The bands were seen before the pre-registration was written; a pre-registration cannot make an
observation on the same data into a prediction. What the confirmatory test does is kill the
confounds, which is worth doing because three claims in this project were lost to exactly one:
a grouping variable coupled to a size, with the statistic reading the size. Here the
grouping variable $\omega$ is collinear with both zero-divisor density and $n$ itself. The
size confound is absent: a blind test reads $\rho(n, \mathrm{Gini}_{\text{add}}) = -0.185$, the
random-orthogonal control has no band structure at all, and a $10{,}000$-shuffle permutation of
$\omega$ puts the minimum adjacent gap at $p = 1.0 \times 10^{-4}$. At matched zero-divisor
density, $10$ discordant pairs favour $\omega$ $10$ out of $10$ (sign test
$p = 1.95 \times 10^{-3}$; the pairs share moduli and are not independent, so the fraction is
the statistic and the $p$ is descriptive).

$\omega$ does not displace zero-divisor density, and we do not crown it. For
predictors $X, Z$ and response $Y$ the partial Spearman correlation
\begin{equation}
  \label{eq:partial}
  \rho(X, Y \mid Z) \;=\;
  \frac{\rho_{XY} - \rho_{XZ}\,\rho_{ZY}}{\sqrt{(1 - \rho_{XZ}^{2})(1 - \rho_{ZY}^{2})}}
\end{equation}
asks what $X$ explains that $Z$ does not. Both survive it:
$\rho(\omega, G \mid \mathrm{zdd}) = +0.804$ and
$\rho(\mathrm{zdd}, G \mid \omega) = +0.578$, so neither is mediated by the other. We decline to
read an ordering off that gap. Ranks here must be \emph{midranks}. $\omega$ takes three distinct
values across the $23$ moduli, and under the ordinal ranking we used until 2026-09-23 the two
partials read $+0.627$ and $+0.735$. Those are the same two numbers in the opposite order,
because within an $\omega$ band an ordinal rank is the ordering of $n$. A quantity whose sign of
difference depends on a tie convention is not a ranking. The band structure was also seen before
the pre-registration, the matched pairs share moduli and are not independent, and the designed
falsifier at $n = 128$ was ambiguous. The defensible statement is that $\omega(n)$ and
zero-divisor density are two partially independent predictors of additive sparsity and this
data cannot rank them.

The designed falsifier fired and neither predictor was accurate. $n = 128 = 2^{7}$ is
a prime power ($\omega = 1$) with half its residues zero divisors, the one place the two
predictors disagree, and the two were pre-registered as disjoint intervals: $\omega$ predicted
$0.017$--$0.184$ and zero-divisor density predicted $0.434$--$0.589$. It read $\mathbf{0.261}$,
between both. We report this as ambiguous. What it settles is smaller than what it was meant
to settle. $128$ lands on $\omega$'s side and extends the $\omega = 1$ band to
$0.015$--$0.261$, while the $\omega = 1$ and $\omega = 2$ bands still do not overlap
($0.261$ against $0.394$). It also favours $\omega$ in all three of the matched pairs it
appears in.

% STE descriptive
\begin{figure}[t]
  \centering
  \includegraphics[width=\textwidth]{../figures/omega_bands.pdf}
  \caption{\textbf{Additive sparsity by $\omega(n)$, with its confound and its falsifier.}
    The panels give the bands, the zero-divisor-density confound, the flat random-orthogonal
    size control, and the $n = 2^{7}$ prediction drawn \emph{beside} its outcome. That
    outcome landed between the two predicted intervals, so neither predictor was accurate.
    The figure regenerates from \texttt{src/viz/plots.py::omega\_bands}.}
  \label{fig:omega}
\end{figure}
% STE end

$\omega$ is the mechanistically expected variable rather than a fished one. The CRT-dual
set of Section~\ref{sec:crt-law} is built as one family per maximal prime power, so at
$\omega = 1$ that set is every frequency, and there is nothing to concentrate on. Each further prime adds
a family and hence more concentration. It is also why a prime power's clock must live in the
multiplicative basis: additively, there is nothing there.

\section{Discussion}
\label{sec:discussion}

This section gathers the interpretations the results support, each at the strength its
evidence carries, and names the measurement that would settle it where one is known.

\paragraph{One mechanism, indexed per stratum.} Every measurable stratum carries a diagonal
clock on its own local group (Section~\ref{sec:strata-result}). Strata that share a local
group select exactly the same characters of it in $158$ of $160$ pre-registered pairs
(Section~\ref{sec:strata-sameness}). At $n = 113$ and $n = 121$, deleting the key characters
from the embedding returns the network to chance (Section~\ref{sec:causal-internal}). The
reading these support is that the network runs the same character-based clock on each
stratum, on that stratum's own local group, and not a private sub-circuit per stratum. The
sameness claim ships without a failed-run control, because the project has only two scorable
failed-run pairs and both come from one model. The same statistic on enough models that never
generalised would complete it.

\paragraph{What the nilpotents do.} \citet{chen2026multiplication} describe non-regular
$\mathcal{J}$-classes as breaking the reduction to local group characters. At every
non-square-free modulus we measured, the reduction does not break. A stratum collapses onto a
smaller local group, $G_{m}$ of Theorem~\ref{thm:stratum}, and the clock runs there
(Sections~\ref{sec:clock-primepowers} and \ref{sec:strata-result}). The coordinate that makes
this measurable was first sighted at one seed (1 seed), and the five-seed measurement of
Section~\ref{sec:strata} supersedes that sighting. The ablation reaches nilpotency depth $7$
at $n = 2^{7}$, and depth beyond that is untested (Section~\ref{sec:causal-128}).

\paragraph{Where the additive structure comes from.} The key additive frequencies of a
grokked model are exactly $P(n)$, one family per maximal prime power
(Section~\ref{sec:crt-law}). The $\mathcal{J}$-class indicators alone concentrate their
energy on the duals of every divisor, and select $P(n)$ at only $2$ of the $10$ composite
moduli tested. The selection of the maximal-prime-power subfamily
is therefore the trained network's, since the algebra alone does not make it. Why the network
makes that selection is not addressed here.

\paragraph{Relation to other mechanistic accounts.} On the unit stratum of a prime the
mechanism is \citet{nanda2023progress}'s clock in discrete-log coordinates, and our engine
reproduces their addition result (Section~\ref{sec:calib}). For $a \cdot b \bmod n$ the two
levels of universality that \citet{chughtai2023toy} describe are algebraic: the family is the
character clock, and the index set is fixed by the stratum's own local group. The top logit
components are $a+b$ terms in discrete-log coordinates at five seeds
(Section~\ref{sec:clock-circuit}), at the one hyperparameter setting we train.
\citet{zhong2023clock} show that other settings can yield other algorithms, and we do not
test them. The basis dependence of \citet{nguyen2026dlogclock} replicates, and the agreement
is specific to float32 (Section~\ref{sec:calib}). The logit ablation survives the three
initialisation conditions of \citet{kunin2024getrich} (Section~\ref{sec:causal-init}).
Section~\ref{sec:limits} names the next test, the same measurement on non-commutative finite
groups.

\paragraph{Precision and absolute numbers.} Training precision changes the measured sparsity
in PyTorch: the same code reads $0.5791$ at float32 and $0.7728$ at float64
(Section~\ref{sec:impl}). A published Gini is therefore comparable to another only at a named
precision, and every absolute sparsity number this paper reports from PyTorch is a float32 measurement.
Comparisons between moduli, and comparisons within one arm, are unaffected. Why our engine
reaches the float64 answer at either dtype is open, with no surviving candidate, and an
operator-by-operator differential trace of one training step at both dtypes would settle it.

\paragraph{What the negatives say about progress measures.} Two statistics that look like
progress measures fire on runs that never generalised. G1, the criterion our pre-registration
named primary, passes on $20$ of $20$ failed stratum-measurements
(Section~\ref{sec:strata-falsifier}). The early CRT enrichment ramp rises monotonically in
failed runs and ends above the band grokked runs occupy early (Section~\ref{sec:crt-ramp}). A
statistic that fires on failures is not evidence, whatever it does on successes. The
procedure we draw from both is to score a candidate statistic on a failed-run control before
proposing it.

\paragraph{Why the shape of this result may travel further than the task.} What the
measurements describe is a network that partitions its input space along a structure the
\emph{data} has, runs one mechanism on every part, and indexes that mechanism by each
part's own symmetry group. Nothing in that sentence mentions modular arithmetic. It is the
pattern an architecture exhibits whenever it routes inputs to sub-computations, as a mixture
of experts does when it chooses among them, or a network whose layers are built to respect a
group action. In those settings the partition is imposed by the designer, where here it was
recovered from a trained model that was never told the algebra of $n$. The instrument transfers with
it: a stratification anyone can compute from the task, a per-stratum statistic, a criterion
fixed before the numbers exist, and a control that has the stratification's size structure
without its content. That last element is what killed three of our own claims
(Appendix~\ref{app:retractions}). We make no claim that a larger model does this; we claim
that this is a testable thing to ask of one.

\section{Limitations}
\label{sec:limits}

% STE descriptive
\paragraph{The early enrichment ramp is not a progress measure.} The CRT enrichment rises
monotonically from early training in grokked runs. It does the same in runs that never
generalise, and it ends higher than grokked runs reach early (Section~\ref{sec:crt-ramp}).
Neither the ramp's presence, its monotonicity nor its permutation $p$ distinguishes a model
that found the circuit from one that memorised its training split. Only the final enrichment
magnitude does. To settle whether any early signal is diagnostic, a statistic must be scored on
a failed-run control before it is proposed. We now apply that procedure, and we did not apply
it when we first measured the ramp.

\paragraph{The engine's dtype-robustness is unexplained.} Section~\ref{sec:impl} shows that
the engine--PyTorch sparsity gap is precision that acts on PyTorch. It leaves open why our
engine reaches the float64 answer at either dtype. We eliminated three candidates under a
pre-registration with both controls: residual type promotion, matmul accumulation order and
softmax formulation. No candidate survives. An operator-by-operator differential trace of a
single training step at both dtypes would settle it.

\paragraph{No algebraic property predicts grokking time.} Grokking time varies by a factor of
$3.8$ across seeds at one modulus. Two claims of ours about how it orders across moduli were
retracted. One was lost to seed variance, and one to a normalisation that manufactured the order
it reported. Nothing in Table~\ref{tab:moduli} predicts it in this data. To establish
that a property does would need many more seeds per modulus than five, at a cost we could
not meet.

\paragraph{$\omega(n)$ is confirmatory, not a discovery.} The additive-sparsity bands were seen
before the pre-registration that tested them, on the same data. A pre-registration cannot
convert that into a prediction. The confounds are dead and $\omega$ does not displace
zero-divisor density, and this data cannot rank the two (Section~\ref{sec:crt-omega}). Its
designed falsifier at $n = 2^{7}$ landed between the two predicted intervals, and neither
predictor was accurate. To rank the two requires moduli chosen to separate them, trained
after the prediction is fixed.

\paragraph{Three claims rest on a single seed.} The torsor coordinate as first sighted, the
vestigial embedding frequency, and the full mechanistic readout of the addition control are
each one seed. They are marked as such wherever they appear. The first is superseded by a
five-seed result, and nothing rests on it. The other two are reported as observations. The four
additional seeds the runs already support would settle each.

\paragraph{Post-grok instability is observed and not explained.} Under weight decay $1.0$,
accuracy excursions below $0.90$ after grokking are common. There are $19$ across $11$ long
runs, all exactly one logging interval wide. Their rate rises with horizon, $0.33\%$ of
samples at $120{,}000$ steps against $0.09\%$ at $40{,}000$. We previously read one such
excursion, which touched a run's final logged sample, as evidence that a network can lose a
learned circuit. It is an artefact of censored data, and the claim is retracted. What causes the
excursions is not addressed here.

A finer logging interval through one excursion would show whether the circuit or the readout
is perturbed.

\paragraph{One test was underpowered.} A search for a specific rare event returned $0$ of
$453$, where the no-effect expectation was $0.68$. The pre-registered probability of zero
events under the effect was also about $0.51$. So the test distinguishes nothing, and its null
is not reported as a result. Only a longer training horizon can expect enough events to decide
it.

\paragraph{Small models, one task family.} Every result comes from a single-layer
transformer, $d_{\text{model}} = 128$, on $a \cdot b \bmod n$ for $n \le 169$. The weight decay
has one value, and so does the training fraction outside the sweep of
Section~\ref{sec:clock-grok}. We do not show that the stratified mechanism is what larger
models or other algebraic tasks do. \citet{zhong2023clock} establish that architectural and
hyperparameter choices can yield different algorithms on a closely related task. The natural
next test is the same measurement on non-commutative finite groups. There,
\citet{chughtai2023toy}'s representation-theoretic account gives an independent prediction of
what the strata should be.

\paragraph{The earliest three sweeps carry no provenance stamp, and their code version is
recovered by argument rather than by record.} We added the provenance stamp after the first
three sweeps ran. So $37$ artifacts carry no stamp at all and $31$ carry the literal string
\texttt{unknown} in place of a commit. Every later sweep, and every run of the from-scratch
engine, carries a real commit. The affected artifacts are behind
Sections~\ref{sec:causal}, \ref{sec:strata} and \ref{sec:crt}, which is most of what this
paper claims. So we state the recovery explicitly rather than leave it implicit.

For the third sweep, the training script has \emph{exactly one} version in the repository's
history. So its code version is determined, and the \texttt{unknown} string records a stamp
failure rather than an ambiguity. For the second sweep, the script has three versions, and we
check their differences mechanically. The model's forward computation and its
non-instrumented return are byte-identical across all three. The only change is the tuple
returned when activations are requested, and the training and evaluation call sites never
request it.

The first sweep, a six-run scout, ran \emph{before} the repository's first commit, so its
history cannot certify what ran. Instead, the source that the compute platform stored for the
version behind those artifacts is byte-identical (equal SHA-256) to the one committed script.
The training path is therefore determined in all three cases.

What remains true is that this is a reconstruction and not a stamp. It rests on a complete
repository history and, for the scout, on the platform's stored record. Every later artifact
rests on a commit recorded at the moment it was written.

The engine's own first run, the addition gate of Section~\ref{sec:calib}, also predates the
repository and was first saved without a stamp. The engine is deterministic and full-batch,
so we re-ran it at a recorded commit. It reproduced every saved weight, logit, activation and
logged value bit for bit. So its artifacts now carry that commit by reproduction rather than
argument.
% STE end

\section{Conclusion}
\label{sec:conclusion}

A transformer trained on $a \cdot b \bmod n$ partitions its input space into algebraic strata
and runs the same character-based clock on each one, on that stratum's own local group. What
the algebra of $n$ changes is not the mechanism but how many strata there are and how they
nest. This supplies the two limitations \citet{chen2026multiplication} state for their own
square-free, correlational result. The reduction to local characters does not break at
non-square-free moduli but relocates to a smaller group. And the frequencies the mechanism
uses are load-bearing rather than merely present: at all $157$ stratum-measurements, in
$28$ of $29$ further single-modulus runs with the exception named in
Section~\ref{sec:causal}, and at all six prime powers. That second answer is an ablation
of the saved logit spectrum, which is the first of the two instruments they ask for. It
bounds what the output depends on, not which component computes it. At $n = 113$ and at
the non-square-free $n = 121$ we also supply the second: deleting those characters from the
embedding and running the network forward takes unit-grid accuracy to $0.0120$ or below
(Section~\ref{sec:causal-internal}). That holds at two moduli, which do not extend to the
other twenty-one.

What it does not supply is a progress measure. The clearest early signal we found rises
monotonically in runs that never generalise, and the criterion our own pre-registration named
primary passes on every failed measurement we scored it against. Both are reported here,
because a statistic that fires on failures is not evidence whatever it does on successes.


\bibliographystyle{tmlr}
\bibliography{refs}

\appendix

\section{The claim ledger}
\label{app:ledger}

Every numbered claim this project made, with its verification status as recorded at the time
of writing. Status is copied from the project's ledger rather than re-judged here. Claims
marked ``(1 seed)'' carry that qualifier in every sentence that makes them. Retracted
claims are in Appendix~\ref{app:retractions} and are not asserted anywhere in the body.

\begin{center}
\footnotesize
\begin{tabular}{l p{0.27\textwidth} p{0.34\textwidth} p{0.16\textwidth}}
  \toprule
  & claim & evidence & status \\
  \midrule
  C1 & algebra table is correct & \texttt{algebra.py::\_selfcheck}, primitive roots by explicit search & verified, self-checking \\
  C2 & non-regular $\mathcal{J}$-class counts $0,0,1,2,8$ & \texttt{j\_structure}; reproduces Prop.\ D.13 on their moduli & verified \\
  C3 & every modulus groks & $23$ moduli, up to $5$ seeds & verified \\
  C4 & $n=113$ replicates prior spectral statistics & $5$ seeds; key-frequency-count clause \textbf{withdrawn} & verified, 5 seeds \\
  C5 & analysis pipeline is correct & run on a third-party published checkpoint & verified externally \\
  C6 & clock survives prime powers & five cyclic prime powers & verified \\
  C7 & non-regular classes carry character structure & superseded by C33 & verified, \textbf{1 seed} \\
  C8 & key additive frequencies are CRT duals & $12$ moduli, $3$--$5$ seeds & verified \\
  C9 & $\mathrm{Gini}_{\text{add}}$ tracks zero-divisor density & $18$ moduli & verified \\
  C10 & from-scratch engine is correct & forward $4.4$e--$16$, gradients $2.68$e--$15$ & verified \\
  C11 & from-scratch engine groks & acceptance test & verified \\
  C12 & softmax collapse occurs on our engine & stamped artifact, $6.2$ orders & verified \\
  C13 & StableMax mitigates C12 & mechanism, not the earlier $10\times$ figure & verified (mechanism) \\
  C14 & critical dataset size is real here & $3$ seeds $\times$ $3$ moduli & verified \\
  C15 & engine reproduces prior addition result & all $8$ criteria & pass \\
  C17 & a frequency can be present yet vestigial & & verified, \textbf{1 seed} \\
  C18 & the addition circuit is a clean Clock & & verified, \textbf{1 seed} \\
  C21 & the circuit uses the embedding's key frequencies & $5$ seeds & verified \\
  C22 & clock frequencies are causally responsible at prime powers & $p<0.01$ in $5/5$ & verified, 5 seeds \\
  C23 & causal where there is no discrete log & $28/29$ analysable runs & verified \\
  C24 & $40$k steps is not converged at four moduli & PR retired at composite moduli & verified \\
  C25 & readout is invariant to the primitive root & exact to $10^{-10}$ & verified \\
  C26 & single-checkpoint Gini is noisy, sd $0.02$--$0.044$ & & verified \\
  C28 & the early CRT signal is CRT-specific & only magnitude discriminates & verified (specificity) \\
  C29 & all four moduli grok on the engine & $3$ seeds, float64 & verified \\
  C31 & the clock is in the neurons & $5$ seeds & verified \\
  C32 & causality survives all three initialisations & $24/24$ & verified \\
  C33 & every stratum runs a clock on its own local group & $32$ strata, $157$ measurements & verified, 5 seeds \\
  C34 & precision changes measured sparsity \emph{in PyTorch} & the $2\times2$ & verified \\
  C35 & $n=113$ is not special & three primes & verified \\
  C36 & the clock is causal at $2^{7}$ & $\hat{p}=10^{-4}$ in $2/2$; $2.77$--$476\times$ & verified \\
  C37 & $\omega(n)$ separates additive sparsity at matched zdd & $23$ moduli & \textbf{confirmatory} \\
  C38 & the same clock across strata, not one per stratum & $158/160$ pairs, null $10^{-4}$ & verified, 5 seeds$^{\dagger}$ \\
  C39 & key characters are what the network computes with & $3/3$ seeds, $n=113$, weight-space & verified \\
  C40 & the same at a non-square-free modulus, non-unit rows intact & $3/3$ seeds, $n=121$, weight-space & verified \\
  \bottomrule
\end{tabular}
\end{center}

\noindent
$^{\dagger}$C38 ships without a working failed-run control: the only two scorable
failed-run pairs in the project come from one model at test accuracy $0.7502$, and its own null
reads $p = 1.0$. Section~\ref{sec:strata-sameness} states this where the claim is made.

\paragraph{On the word \emph{causal} in this table.} These rows are the project's claim ledger,
quoted in the wording each claim was registered under. Throughout, \emph{causal} and
\emph{causally responsible} mean the logit-spectrum ablation of
Section~\ref{sec:methods-causal} (restricted and excluded losses against a permutation null),
and never weight ablation or activation patching, neither of which this paper performs.

\section{Retractions}
\label{app:retractions}

Five claims were made and withdrawn. Each is recorded with what killed it and the control that
would have caught it earlier. No paper in the corpus this one joins does this; we do it because
it is what makes the surviving claims checkable.

Three of the five died to the same shape, and it is worth stating once in symbols. A grouping
variable $X$ is compared against a response $Y$, but $X$ is structurally coupled to a size $S$
that $Y$ also tracks:
\begin{equation}
  \label{eq:confound}
  X \longrightarrow S \longrightarrow Y,
  \qquad\text{so } \rho(X, Y) \text{ is large while } \rho(X, Y \mid S) \text{ is not,}
\end{equation}
with $\rho(\cdot \mid \cdot)$ the partial correlation of \eqref{eq:partial}. The cheap general
defence, which caught all three, is to run the statistic on a control that has the size
structure and not the effect: shuffled labels, flat variance, a permuted grouping. If the
control reproduces the result, there is no result.

\paragraph{C19: $\varphi(n)$-normalised grokking time separates cyclic from non-cyclic
moduli.} \emph{Killed by:} a size confound in the normalisation itself. Raw grokking time
\emph{decreases} with $\varphi$ ($\rho = -0.665$), so dividing by $\varphi$ over-corrects: with
$S = \varphi(n)$ in \eqref{eq:confound}, $\rho(-\varphi, \text{steps}/\varphi) = +0.870$ over
$15$ moduli, an ordering manufactured entirely by the normalisation. The
cyclic moduli in the original sample were simply the three largest $\varphi$. \emph{The control
that would have caught it:} correlating the normalised statistic with the quantity it was
normalised by, before interpreting it.

\paragraph{C20: activation variance is graded by $\mathcal{J}$-class depth.} \emph{Killed
by:} a size confound between depth and block cell count, coupled through
$|J_{d}| = \varphi(n/d)$ in \eqref{eq:jclass}. $\rho(\text{cells},
\text{variance})$ runs $+0.856$ to $+0.949$, and a flat-variance \emph{noise} control scores
$\rho = +0.942$. \emph{The control:} measuring every block at an equal subsampled cell count,
and dropping blocks that cannot supply it rather than reporting them as zero.

\paragraph{C7b: regular $\mathcal{J}$-classes carry more structure than non-regular ones.}
\emph{Killed by:} a size confound between regularity and class size: $J_{1}$ is always
regular and, by \eqref{eq:jclass}, always the largest class. Size-matched within one model at $n = 120$, the only place such
pairs exist, the two overlap. \emph{The control:} the size-matched comparison, which is what
retired it.

\paragraph{C27: a network can lose a learned circuit.} \emph{Killed by:} right-censoring. The
claim was one run read at its final logged sample. Across $11$ long runs there are $19$
post-grok excursions below accuracy $0.90$, every one exactly one logging interval wide, $18$
of which recover. The nineteenth merely begins at the last step at which an observation
exists, and that run had recovered from a deeper excursion $4{,}600$ steps earlier. \emph{The control:}
classifying a run by a window rather than by its last row, which is now done everywhere.

\paragraph{C30: training precision changes measured sparsity.} \emph{Killed by:} a controlled
A/B that varied dtype inside our engine alone and returned the wrong sign, effect size
$-0.07$. It is superseded by C34, which is the same phenomenon as an \emph{interaction}:
precision is decisive in PyTorch and inert in our engine. C30's main-effect sentence is never
revived. \emph{The control:} running the variable. The claim had a $7/7$ correlation and a
mechanism that tied up three loose ends, which is the most persuasive shape a wrong claim can
take, and no amount of care in the statistics would have rescued it.

\section{Pre-registrations}
\label{app:prereg}

Every experiment whose criteria could be fixed in advance has a pre-registration committed
before the run, and the git timestamp is the evidence the criteria predated the result. Each is
named in the body beside the result it governs, with its commit. Two exceptions are stated
where they occur: the first gate's criteria, two of which were edited after failing
(Section~\ref{sec:calib}), and the $\omega(n)$ analysis, which is confirmatory because the
bands were seen before the pre-registration was written.

\section{Per-modulus and per-stratum tables}
\label{app:tables}

The full per-stratum table underlying Section~\ref{sec:strata}, and per-modulus spectral
statistics for all $23$ moduli, regenerate from the repository.

\section{Instruments}
\label{app:instruments}

The from-scratch autograd engine, its finite-difference test suite, the memory guards, and the
self-check suite that gates every claim-carrying script. One item is worth recording as a
limitation of its own: the engine's reported speedup is verified only for numerics, to
$2.68 \times 10^{-15}$; the wall-clock figure is unstamped and machine-dependent.

\section{Experimental details and reproduction}

\subsection{Code and data availability}
\label{sec:availability}

Every number in this paper regenerates from the code released with it (the supplementary
material during review). The release holds the autograd engine, the transformer, the analysis
toolkit, the kernels that produced each sweep, and \texttt{scripts/reproduce.sh}. That script
re-runs every analysis from the saved artifacts, with pinned dependencies and fixed seeds. It
does not train a model (Section~\ref{sec:repro-procedure}). Saved artifacts carry a provenance stamp
recording the git SHA, the dtype, the account whose compute produced them and the exact
configuration. \texttt{test\_paper\_numbers.py} re-derives every load-bearing number in this
manuscript, captions included, from those artifacts rather than from the prose, and it runs as
part of the self-check suite. Two limitations of that guarantee are stated in
Section~\ref{sec:limits}. $68$ artifacts from the three earliest sweeps carry no usable SHA,
and their code version is recovered by argument rather than by stamp. Figures are regenerated
rather than stored, so a figure is only as reproducible as the artifact beneath it.

\subsection{The sweeps}
\label{sec:sweeps}

% STE descriptive
Tables~\ref{tab:sweeps} and~\ref{tab:sweep-files} list every sweep that supplies a result. The
first table gives the configuration of each sweep. The second table gives the script that
trained it, the directory of its artifacts, the script that analyses it and its compute.
\texttt{scripts/sweep\_table.py} writes both tables from the artifacts and the run logs.
\texttt{test\_paper\_numbers.py} checks that the tables here are equal to its output.

\begin{table}[ht]
  \centering
  \footnotesize
  \caption{\textbf{The configuration of each sweep.} ``Engine'' is our from-scratch
    implementation and ``PyTorch'' is the reference implementation. Every sweep uses AdamW
    with learning rate $10^{-3}$ and weight decay $1.0$, full-batch. A moduli cell with more
    than four moduli gives their count and range. Table~\ref{tab:moduli} lists each modulus.}
  \label{tab:sweeps}
  \begin{tabular}{llllllll}
    \toprule
    sweep & code & precision & moduli & seeds & steps & train fraction & platform \\
    \midrule
    k01 & PyTorch & float32 & $6$, $113$ to $165$ & $0$ & $25{,}000$ & $0.3$ & Kaggle T4 \\
    k02 & PyTorch & float32 & $6$, $113$ to $165$ & $0$--$4$ & $40{,}000$ & $0.3$ & Kaggle T4 \\
    k03 & PyTorch & float32 & $6$, $113$ to $165$ & $0$--$4$ & $40{,}000$ & $0.3$ & Kaggle T4 \\
    k04 & PyTorch & float32 & $12$, $49$ to $169$ & $0$--$2$ & $40{,}000$ & $0.3$ & Kaggle T4 \\
    k05 & PyTorch & float32 & $49$, $54$, $63$ & $0$--$2$ & $40{,}000$ & $0.3$, $0.5$, $0.8$ & Kaggle T4 \\
    k06 & PyTorch & float32 & $119$, $120$, $125$, $165$ & $0$--$2$ & $120{,}000$ & $0.3$ & Kaggle T4 \\
    k07 & PyTorch & float32 & $113$ & $0$--$4$ & $40{,}000$ & $0.3$ & Kaggle T4 \\
    k08 & PyTorch & float32 & $113$ & $0$--$7$ & $40{,}000$ & $0.3$ & Kaggle T4 \\
    k09 & PyTorch & float32 & $5$, $91$ to $131$ & $0$--$2$ & $40{,}000$ & $0.3$ & Kaggle T4 \\
    N7 & engine & float64 & $113$, $119$, $121$, $125$ & $0$--$2$ & $40{,}000$ & $0.3$ & CPU \\
    N9 & engine & float32 & $113$ & $0$--$2$ & $40{,}000$ & $0.3$ & CPU \\
    O18 & PyTorch & float32 & $113$ & $0$--$2$ & $40{,}000$ & $0.3$ & CPU \\
    O18, f64 & PyTorch & float64 & $113$ & $0$--$2$ & $40{,}000$ & $0.3$ & CPU \\
    O18-T & engine & float64 & $113$ & $0$ & $40{,}000$ & $0.3$ & CPU \\
    I1 & engine & float64 & $113$ & $0$--$2$ & $40{,}000$ & $0.3$ & CPU \\
    I1b & engine & float64 & $121$ & $0$--$2$ & $40{,}000$ & $0.3$ & CPU \\
    Gate 1 & engine & float64 & $113$ & $0$ & $25{,}000$ & $0.3$ & CPU \\
    \bottomrule
  \end{tabular}
\end{table}

\begin{table}[ht]
  \centering
  \footnotesize
  \caption{\textbf{The files and the compute of each sweep.} N9 runs \texttt{run\_n7.py}
    with \texttt{ENGINE\_DTYPE=float32}. I1 and I1b train with \texttt{run\_n7.py} and then
    measure with \texttt{run\_intervention.py}. The hours of a Kaggle sweep are the wall time
    of its T4 session. The hours of a CPU sweep are the sum of the wall times of its runs,
    and some of those runs were in parallel.}
  \label{tab:sweep-files}
  \begin{tabular}{lllll}
    \toprule
    sweep & training script & artifact directory & analysis script & hours \\
    \midrule
    k01 & \texttt{kernels/k01\_scout/run.py} & \texttt{results/k01\_scout} & \texttt{analyze\_scout.py} & $0.6$ \\
    k02 & \texttt{kernels/k02\_grid/run.py} & \texttt{results/k02\_grid} & \texttt{analyze\_k02.py} & $3.7$ \\
    k03 & \texttt{kernels/k03\_grid\_acts/run.py} & \texttt{results/k03\_grid\_acts} & \texttt{analyze\_k03.py} & $3.8$ \\
    k04 & \texttt{kernels/k04\_extended/run.py} & \texttt{results/k04\_extended} & \texttt{analyze\_k04.py} & $3.2$ \\
    k05 & \texttt{kernels/k05\_lowdata/run.py} & \texttt{results/k05\_lowdata} & \texttt{analyze\_k05.py} & $1.3$ \\
    k06 & \texttt{kernels/k06\_horizon/run.py} & \texttt{results/k06\_horizon} & \texttt{analyze\_k06.py} & $4.7$ \\
    k07 & \texttt{kernels/k07\_arma\_seeds/run.py} & \texttt{results/k07\_arma\_seeds} & \texttt{analyze\_k07.py} & $0.5$ \\
    k08 & \texttt{kernels/k08\_init/run.py} & \texttt{results/k08\_init} & \texttt{analyze\_k08.py} & $2.2$ \\
    k09 & \texttt{kernels/k09\_primes\_zdd/run.py} & \texttt{results/k09\_primes\_zdd} & \texttt{analyze\_k09.py} & $1.6$ \\
    N7 & \texttt{run\_n7.py} & \texttt{results/n7\_engine} & \texttt{analyze\_n7.py} & $68.5$ \\
    N9 & \texttt{run\_n7.py} & \texttt{results/n9\_f32} & \texttt{analyze\_n9.py} & $9.6$ \\
    O18 & \texttt{run\_o18\_cpu.py} & \texttt{results/o18\_cpu} & \texttt{analyze\_n9.py} & $1.8$ \\
    O18, f64 & \texttt{run\_o18\_cpu.py -{}-f64} & \texttt{results/o18\_f64} & \texttt{analyze\_n9.py} & $3.2$ \\
    O18-T & \texttt{run\_o18\_transplant.py} & \texttt{results/o18\_transplant} & \texttt{run\_o18\_transplant.py} & $5.3$ \\
    I1 & \texttt{run\_n7.py} & \texttt{results/i1\_internal} & \texttt{run\_intervention.py} & $20.7$ \\
    I1b & \texttt{run\_n7.py} & \texttt{results/i1\_internal} & \texttt{run\_intervention.py} & $24.3$ \\
    Gate 1 & \texttt{run\_gate1.py} & \texttt{results/gate1} & \texttt{test\_gate1.py} & $3.1$ \\
    \bottomrule
  \end{tabular}
\end{table}
% STE end

\subsection{Reproduction procedure}
\label{sec:repro-procedure}

% STE procedural
To reproduce the numbers and the figures of this paper, do these steps:
\begin{enumerate}
  \item Unpack the supplementary archive.
  \item Install Python 3.12.13.
  \item Install the pinned dependencies with \texttt{pip install -r requirements.txt}.
  \item Put the saved artifacts in \texttt{results/} and the run logs in \texttt{logs/}.
  \item From the top directory, run \texttt{bash scripts/reproduce.sh}.
\end{enumerate}
% STE end

% STE descriptive
The script does not train a model. It runs the self-checks and the analysis scripts of
Table~\ref{tab:sweep-files}. The self-checks test Proposition~\ref{prop:torsor} and
Theorem~\ref{thm:stratum} by exhaustion at every modulus up to $200$. The script also runs
\texttt{test\_paper\_numbers.py}, which calculates the numbers in the text, the tables and the
captions again from the artifacts. It fails on a number that it does not calculate, unless
\texttt{paper/number\_ledger.json} records that number with a reason.

\texttt{scripts/render\_all.py} then draws every figure again.
\texttt{scripts/modulus\_table.py} writes Table~\ref{tab:moduli}, and
\texttt{scripts/sweep\_table.py} writes Tables~\ref{tab:sweeps} and~\ref{tab:sweep-files}. The
script needs no GPU. Its longest step is \texttt{analyze\_gate2.py} on the k03 sweep, which
took $8.5$ hours on our CPU.
% STE end

% STE procedural
To train a sweep again, do these steps:
\begin{enumerate}
  \item Run \texttt{scripts/sweep\_table.py -{}-commands}.
  \item Find the sweep. The output gives its git SHA and its exact command.
  \item Run that command from the top directory.
\end{enumerate}
% STE end

% STE descriptive
The self-check of \texttt{sweep\_table.py} makes sure that each command issues the arguments
that the artifacts of its sweep record. It also makes sure that the training code at the
current commit is the same as at that SHA. Where the code is different, the self-check makes
sure that the difference cannot change a result. A Kaggle
sweep runs on a GPU without deterministic settings. Thus its second run is not guaranteed to be
bit-identical to the first (Section~\ref{sec:impl}).
% STE end


\end{document}
